\chapter{命题逻辑}
逻辑，是英文单词 logic 的音译.
逻辑学是研究思维形式、思维方法及思维规律尤其是推理的学科.
早在两千多年前，古希腊逻辑学家亚里士多德创立了形式逻辑（formal logic），
详细讨论概念（词项）、判断（命题）和各种形式的推理，
研究逻辑基本规律等内容.

德国数学家、哲学家莱布尼茨首先提出用数学方法研究逻辑，
即建立一套表意符号体系，在符号之间进行形式推理，
由此创立了数理逻辑（mathematical logic）.

现在，除了传统的数理逻辑（内容包括逻辑演算、公理化集合论、模型论、递归论和证明论）以外，
还出现了各种各样的应用逻辑，诸如
多值逻辑、模态逻辑、归纳逻辑、时序逻辑、动态逻辑、模糊逻辑、
非单调逻辑、默认逻辑、数字逻辑、电路逻辑、算法逻辑及程序逻辑等，
这些都与计算机科学密切相关.

本章介绍\DefineConcept{命题逻辑}（propositional logic）.
%@see: https://www.tutorialspoint.com/discrete_mathematics/rules_of_inference.htm

\section{命题，形式语言}
\subsection{命题的概念}
\DefineConcept{命题}（proposition），
是指对确定的对象进行判断的陈述句.
这可以从两个方面去理解：\begin{enumerate}
	\item 命题必须是一个完整的句子，具有必需的语法结构（起码要有主语、谓语、宾语）；
	\item 命题必须具有真假意义，研究者应该能够判断命题是否符合客观实际或是否合理，
	这就要求命题不能是除了陈述句以外的其他句式（例如疑问句、感叹句）.
\end{enumerate}

下面几句话都是命题：\begin{itemize}
	\item 太阳总从东方升起.
	\item 只有在冬天才下雪.
	\item 水是液体.
	\item 冰是固体.
	\item 辽宁舰是中国的第一艘航空母舰.
	\item 我喜欢智能手机盒平板电脑.
	\item 小李和小王是同学.
	\item 你只有刻苦学习，才能取得好成绩.
	\item 火星上有生物\footnote{
		虽然我们至今尚不知道火星上是否存在生物，
		但我们相信在将来某个时刻一定会知道的，
		因此“火星上有生物.”是一个命题.
	}.
\end{itemize}

“\(x>3\).”不是命题，因为我们无从得知变量\(x\)的取值，进而无法确定其真假.

“立正！”是命令句，它本身没有对错之分，但命令发出后会有终结反应.

“这朵花真漂亮！”是感叹句，
“你要我的手机号码是想给我充话费？”是疑问句，
都没有真假意义.

命题的\DefineConcept{真值}（truth）就是命题的逻辑取值.
经典逻辑值只有两个：\begin{itemize}
	\item 如果命题\(p\)的判断正确，
	那么我们称“命题\(p\)是\DefineConcept{真的}（true）”，
	或称“命题\(p\)是\DefineConcept{真命题}”，
	或称“命题\(p\)的真值为\(1\)”；
	\item 如果命题\(p\)的判断错误，
	那么我们称“命题\(p\)是\DefineConcept{假的}（false）”，
	或称“命题\(p\)是\DefineConcept{假命题}”，
	或称“命题\(p\)的真值为\(0\)\footnote{
		有的书会把真命题的真值记为T，
		把假命题的真值记为F.
		还有的会把真命题的真值记为\(\top\)，
		把假命题的真值记为\(\bot\).
	}”.
\end{itemize}

\subsection{命题逻辑的形式语言}
对于用汉语（或其他人类语言）表述的命题，
我们可能会因为它含糊不清的语义而无法作出真伪判断.
因此，数学家创造了一种独立于人类语言的、专门用于逻辑演绎的语言，
这就是\DefineConcept{形式语言}（formal language）.

形式语言有三个方面可供我们讨论：\begin{enumerate}
	\item 指定形式语言使用的符号集、字母表（即下面即将介绍的命题常量、命题变元）；
	\item 制定一些规则，用以构造语法正确的有限符号串（即命题公式）；
	\item 指明形式语言与自然语言之间所允许的翻译.
\end{enumerate}

我们把用来表示命题真值的符号\(0\)和\(1\)
称为\DefineConcept{逻辑常量}或\DefineConcept{命题常量}.
相对地，把用来表示命题的小写拉丁字母（例如\(p,q,r,\dotsc\)）
称为\DefineConcept{逻辑变量}或\DefineConcept{命题变元}.
命题变元可以代表任意命题.
从取值的角度看，命题变元既可以取\(1\)也可以取\(0\).

\section{逻辑联结词，命题公式，真值表}
对于复杂的命题，我们可以将其看作是由若干个小命题组合而成的.
就像汉语中有“而且”“但是”等连接句子的虚词一样，
命题逻辑也有扮演相同角色的符号，
这就是\DefineConcept{逻辑联结词}（logical connective）.

\DefineConcept{命题公式}（proposition formula）是
由命题常量、命题变元、逻辑联结词、左圆括号、右圆括号构成的有意义的符号串，
本质上是逻辑函数或逻辑表达式，其取值只可能为\(1\)或\(0\)，
其严格定义可以借助递归定义方式给出.
\begin{definition}
%@see: 《离散数学》（邓辉文） P85 定义3-1
命题公式按下列规则生成：\begin{enumerate}[label={(\arabic*)}]
	\item 命题常量、命题变元是命题公式；
	\item 若\(A\)是命题公式，则\((\neg A)\)是命题公式；
	\item 若\(A\)和\(B\)是命题公式，则\((A \circ B)\)是命题公式，
	其中符号\(\circ\)代表某个逻辑联结词；
	\item 有限次应用(1)(2)(3)这三条规则所得到的符号串是命题公式.
\end{enumerate}
\end{definition}

我们可以按命题公式中是否含有逻辑联结词，将命题公式分为两类：\begin{itemize}
	\item \DefineConcept{简单命题}（simple proposition），
	或称\DefineConcept{原子命题}（atom proposition），
	即不含有逻辑联结词的命题公式；
	\item \DefineConcept{复合命题}（compound proposition），
	即由简单命题和逻辑联结词构成的命题公式.
\end{itemize}

\subsection{命题公式的真值指派、真值表}
给定命题公式\(A\)，对\(A\)中出现的每个命题变元都指定一个真值\(1\)或\(0\)
\footnote{
	目前而言，真值集合是\(\{1,0\}\)，不含其他值.
	但是真值集合也可以是三值的，还可以是\(\mathbb{N}\)或\([0,1]\)等等.
}，
就说“对命题公式\(A\)进行了一种\DefineConcept{真值指派}（assignment）”
或者说“对命题公式\(A\)进行了一个\DefineConcept{解释}（interpretation）”，
而在该指派下可以求出命题公式\(A\)的一个真值.
将\(A\)的所有可能的真值指派以及它在每一个真值指派下的取值列成一张表，
就得到命题公式\(A\)的\DefineConcept{真值表}（truth table）.

\subsection{否定联结词}
设\(p\)表示一个命题，
\(\neg p\)是
“对命题\(p\)的\DefineConcept{否定}（negation）”，
读作“非\(p\)”.

\begin{table}[ht]
%@Mathematica: TableForm[BooleanTable[{p, Not[p]}, {p}] // Boole, TableHeadings -> {None, {"p", "\[Not]p"}}]
	\centering
	\begin{tabular}{|c|p{1.5cm}|}
		\hline
		\(p\) & \(\neg p\) \\
		\hline
		0 & 1 \\
		1 & 0 \\
		\hline
	\end{tabular}
	\caption{否定联结词的真值表}
\end{table}

\(\neg\)是本节采用的唯一一种一元逻辑运算符，
下面讨论的都是二元逻辑运算符.

\subsection{合取联结词}
设\(p,q\)各表示一个命题，
\(p \land q\)是
“对命题\(p\)与\(q\)的\DefineConcept{合取}（conjunction）”.

合取联结词\(\land\)相当于汉语中的
“并且”“与”“和”“以及”“不但……而且……”
“虽然……但是……”“尽管……仍然……”等.

\begin{table}[ht]
%@Mathematica: TableForm[BooleanTable[{p, q, And[p, q]}, {p, q}] // Boole, TableHeadings -> {None, {"p", "q", "p \[And] q"}}]
	\centering
	\begin{tabular}{|*{2}{c|}p{2cm}|}
		\hline
		\(p\) & \(q\) & \(p \land q\) \\
		\hline
		0 & 0 & 0 \\
		0 & 1 & 0 \\
		1 & 0 & 0 \\
		1 & 1 & 1 \\
		\hline
	\end{tabular}
	\caption{合取联结词的真值表}
\end{table}

\subsection{析取联结词}
设\(p,q\)各表示一个命题，
\(p \lor q\)是
“对命题\(p\)与\(q\)的\DefineConcept{析取}（disjunction）”.

析取联结词\(\lor\)相当于汉语中的“或”.

\begin{table}[ht]
%@Mathematica: TableForm[BooleanTable[{p, q, Or[p, q]}, {p, q}] // Boole, TableHeadings -> {None, {"p", "q", "p \[Or] q"}}]
	\centering
	\begin{tabular}{|*{2}{c|}p{2cm}|}
		\hline
		\(p\) & \(q\) & \(p \lor q\) \\
		\hline
		0 & 0 & 0 \\
		0 & 1 & 1 \\
		1 & 0 & 1 \\
		1 & 1 & 1 \\
		\hline
	\end{tabular}
	\caption{析取联结词的真值表}
\end{table}

\subsection{异或联结词}
设\(p,q\)各表示一个命题，
\(p \lxor q\)是
“对命题\(p\)与\(q\)的\DefineConcept{异或}（exclusive or）”.

为了凸出析取联结词与异或联结词的联系与区别，
我们常把对命题\(p\)与\(q\)的析取\(p \lor q\)
称为“对\(p\)与\(q\)的\DefineConcept{可兼或}（inclusive or）”，
把对\(p\)与\(q\)的异或\(p \lxor q\)
称为“对\(p\)与\(q\)的\DefineConcept{不可兼或}”.

\begin{table}[ht]
%@Mathematica: TableForm[BooleanTable[{p, q, Xor[p, q]}, {p, q}] // Boole, TableHeadings -> {None, {"p", "q", "p \[Xor] q"}}]
	\centering
	\begin{tabular}{|*{2}{c|}p{2cm}|}
		\hline
		\(p\) & \(q\) & \(p \lxor q\) \\
		\hline
		0 & 0 & 0 \\
		0 & 1 & 1 \\
		1 & 0 & 1 \\
		1 & 1 & 0 \\
		\hline
	\end{tabular}
	\caption{异或联结词的真值表}
\end{table}

\subsection{条件联结词}
%@see: https://plato.stanford.edu/entries/necessary-sufficient/
设\(p,q\)各表示一个命题，
\(p \limp q\)读作“\(p\)蕴含\(q\)”或“\(p\)条件\(q\)”.

我们把\(p\)称为“\(p \limp q\)的\DefineConcept{前件}（antecedent）”
或“\(p \limp q\)的\DefineConcept{条件}（premise）”，
把\(q\)称为“\(p \limp q\)的\DefineConcept{后件}（consequent）”
或“\(p \limp q\)的\DefineConcept{结论}（conclusion）”.

\begin{table}[ht]
%@Mathematica: TableForm[BooleanTable[{p, q, Implies[p, q]}, {p, q}] // Boole, TableHeadings -> {None, {"p", "q", "p \[Implies] q"}}]
	\centering
	\begin{tabular}{|*{2}{c|}p{2cm}|}
		\hline
		\(p\) & \(q\) & \(p \limp q\) \\
		\hline
		0 & 0 & 1 \\
		0 & 1 & 1 \\
		1 & 0 & 0 \\
		1 & 1 & 1 \\
		\hline
	\end{tabular}
	\caption{条件联结词的真值表}
\end{table}

如果已知命题\(p \limp q\)是真命题，
则称“\(p\)是\(q\)的一个\DefineConcept{充分条件}（sufficient condition）”，
或称“\(q\)是\(p\)的一个\DefineConcept{必要条件}（necessary condition）”.

如果已知\((p \limp q)\)和\((q \limp p)\)都是真命题，
则称“\(p\)是\(q\)的一个\DefineConcept{充分必要条件}”，
或称“\(q\)是\(p\)的一个\DefineConcept{充分必要条件}”.

如果已知\((p \limp q)\)是真命题而\((q \limp p)\)是假命题，
则称“\(p\)是\(q\)的一个\DefineConcept{充分不必要条件}”，
或称“\(q\)是\(p\)的一个\DefineConcept{必要不充分条件}”.

把\(q \limp p\)称为“\(p \limp q\)的\DefineConcept{逆命题}”.
称“\(q \limp p\)和\(p \limp q\)是\DefineConcept{互逆命题}”.

把\(\neg p \limp \neg q\)称为“\(p \limp q\)的\DefineConcept{否命题}”.
称“\(\neg p \limp \neg q\)和\(p \limp q\)是\DefineConcept{互否命题}”.

把\(\neg q \limp \neg p\)称为“\(p \limp q\)的\DefineConcept{逆否命题}”.
称“\(\neg q \limp \neg p\)和\(p \limp q\) \DefineConcept{互为逆否命题}”.

\begin{figure}[htb]
	\centering
	\tikzstyle{prop} = [rectangle, minimum width=3cm, minimum height=2cm, text centered, draw=black, fill=orange!30]
	\tikzstyle{arrow} = [thick,<->,>=stealth]
	\begin{tikzpicture}[node distance=4cm]
		\node (p11) [prop] {\begin{tblr}{c}
			原命题 \\ \(p \limp q\)
		\end{tblr}};
		\node (p12) [prop, right=6cm of p11] {\begin{tblr}{c}
			逆命题 \\ \(q \limp p\)
		\end{tblr}};
		\node (p21) [prop, below of=p11] {\begin{tblr}{c}
			否命题 \\ \(\neg p \limp \neg q\)
		\end{tblr}};
		\node (p22) [prop, below of=p12] {\begin{tblr}{c}
			逆否命题 \\ \(\neg q \limp \neg p\)
		\end{tblr}};

		\begin{scope}[arrow]
			\draw (p11) -- node[anchor=south]{互逆} (p12);
			\draw (p11) -- node[anchor=east]{互否} (p21);
			\draw (p22) -- node[anchor=north]{互逆} (p21);
			\draw (p22) -- node[anchor=west]{互否} (p12);
			\draw (p11) -- node[sloped,near end,below]{互为逆否} (p22);
			\draw (p12) -- node[sloped,near end,above]{互为逆否} (p21);
		\end{scope}
	\end{tikzpicture}
	\caption{}
\end{figure}

\subsection{等价联结词}
设\(p,q\)各表示一个命题，
\(p \liff q\)读作“\(p\)等价\(q\)”或“\(p\)双条件\(q\)”.

\begin{table}[ht]
%@Mathematica: TableForm[BooleanTable[{p, q, Equivalent[p, q]}, {p, q}] // Boole, TableHeadings -> {None, {"p", "q", "p \[Equivalent] q"}}]
	\centering
	\begin{tabular}{|*{2}{c|}p{2cm}|}
		\hline
		\(p\) & \(q\) & \(p \liff q\) \\
		\hline
		0 & 0 & 1 \\
		0 & 1 & 0 \\
		1 & 0 & 0 \\
		1 & 1 & 1 \\
		\hline
	\end{tabular}
	\caption{等价联结词的真值表}
\end{table}

\subsection{与非联结词}
设\(p,q\)各表示一个命题，
\(p \lnand q\)读作“\(p\)与非\(q\)”.

\begin{table}[ht]
%@Mathematica: TableForm[BooleanTable[{p, q, Nand[p, q]}, {p, q}] // Boole, TableHeadings -> {None, {"p", "q", "p \[Nand] q"}}]
	\centering
	\begin{tabular}{|*{2}{c|}p{2cm}|}
		\hline
		\(p\) & \(q\) & \(p \lnand q\) \\
		\hline
		0 & 0 & 1 \\
		0 & 1 & 1 \\
		1 & 0 & 1 \\
		1 & 1 & 0 \\
		\hline
	\end{tabular}
	\caption{与非联结词的真值表}
\end{table}

\subsection{或非联结词}
设\(p,q\)各表示一个命题，
\(p \lnor q\)读作“\(p\)或非\(q\)”.

\begin{table}[ht]
%@Mathematica: TableForm[BooleanTable[{p, q, Nor[p, q]}, {p, q}] // Boole, TableHeadings -> {None, {"p", "q", "p \[Nor] q"}}]
	\centering
	\begin{tabular}{|*{2}{c|}p{2cm}|}
		\hline
		\(p\) & \(q\) & \(p \lnor q\) \\
		\hline
		0 & 0 & 1 \\
		0 & 1 & 0 \\
		1 & 0 & 0 \\
		1 & 1 & 0 \\
		\hline
	\end{tabular}
	\caption{或非联结词的真值表}
\end{table}

\subsection{条件否定联结词}
设\(p,q\)各表示一个命题，
\(p \lnimp q\)读作“\(p\)条件否定\(q\)”.

\begin{table}[ht]
%@Mathematica: TableForm[BooleanTable[{p, q, Not[Implies[p, q]]}, {p, q}] // Boole, TableHeadings -> {None, {"p", "q", "p !\[Implies] q"}}]
	\centering
	\begin{tabular}{|*{2}{c|}p{2cm}|}
		\hline
		\(p\) & \(q\) & \(p \lnimp q\) \\
		\hline
		0 & 0 & 0 \\
		0 & 1 & 0 \\
		1 & 0 & 1 \\
		1 & 1 & 0 \\
		\hline
	\end{tabular}
	\caption{条件否定联结词的真值表}
\end{table}

\subsection{括号的省略}
我们解读形式语言的顺序和汉语、英语等常见人类语言一致.
运算次序由括号、联结词的优先顺序及下面的结合约定共同决定，
不能仅根据从左向右的阅读顺序确定.

严格按照命题公式的定义，就会出现很多的括号.
一方面，这些括号使得命题公式的结构清晰、含义清除；
而另一方面，括号太多给命题公式的阅读和书写带来不便.
因此，为了简化书写，减少命题公式中括号的数目，特作如下一些可以省略括号的约定：\begin{itemize}
	\item 最外层的括号可以省略；
	\item 9个逻辑联结词运算的优先顺序依次为\footnote{
		有的书上只规定了\(\neg,\land,\lor,\limp,\liff\)这5个逻辑联结词的优先顺序.
	}\begin{equation*}
		\neg, \quad
		\land, \quad
		\lor, \quad
		\lxor, \quad
		\limp, \quad
		\liff, \quad
		\lnand, \quad
		\lnor, \quad
		\lnimp,
	\end{equation*}
	符合本约定的有些括号可以不写（详见\cref{definition:命题逻辑.逻辑联结词的优先顺序}）；
	\item 同级运算从右往左依次进行
	\footnote{
		有的书上规定同级运算从左往右依次进行.
		以后我们会证明，
		当改变运算顺序时，
		\(\land,\lor,\lxor,\liff\)的运算结果不变，
		\(\limp,\lnand,\lnor\)的运算结果改变.
	}.
\end{itemize}

\begin{definition}\label{definition:命题逻辑.逻辑联结词的优先顺序}
设\(\phi,\psi,\eta\)都是命题公式.
定义：\begin{gather}
	\phi
	\defiff
	(\phi), \\
	(\neg\phi\land\psi)
	\defiff
	((\neg\phi)\land\psi), \\
	(\phi\land\neg\psi)
	\defiff
	(\phi\land(\neg\psi)), \\
	(\neg\phi\lor\psi)
	\defiff
	((\neg\phi)\lor\psi), \\
	(\phi\lor\neg\psi)
	\defiff
	(\phi\lor(\neg\psi)), \\
	(\neg\phi\limp\psi)
	\defiff
	((\neg\phi)\limp\psi), \\
	(\phi\limp\neg\psi)
	\defiff
	(\phi\limp(\neg\psi)), \\
	(\neg\phi\liff\psi)
	\defiff
	((\neg\phi)\liff\psi), \\
	(\phi\liff\neg\psi)
	\defiff
	(\phi\liff(\neg\psi)), \\
	(\phi\land\psi\land\eta)
	\defiff
	(\phi\land(\psi\land\eta)), \\
	(\phi\land\psi\lor\eta)
	\defiff
	((\phi\land\psi)\lor\eta), \\
	(\phi\lor\psi\lor\eta)
	\defiff
	(\phi\lor(\psi\lor\eta)), \\
	(\phi\lor\psi\land\eta)
	\defiff
	(\phi\lor(\psi\land\eta)), \\
	(\phi\land\psi\limp\eta)
	\defiff
	((\phi\land\psi)\limp\eta), \\
	(\phi\limp\psi\limp\eta)
	\defiff
	(\phi\limp(\psi\limp\eta)), \\
	(\phi\limp\psi\land\eta)
	\defiff
	(\phi\limp(\psi\land\eta)), \\
	(\phi\land\psi\liff\eta)
	\defiff
	((\phi\land\psi)\liff\eta), \\
	(\phi\liff\psi\liff\eta)
	\defiff
	(\phi\liff(\psi\liff\eta)), \\
	(\phi\liff\psi\land\eta)
	\defiff
	(\phi\liff(\psi\land\eta)), \\
	(\phi\lor\psi\limp\eta)
	\defiff
	((\phi\lor\psi)\limp\eta), \\
	(\phi\limp\psi\lor\eta)
	\defiff
	(\phi\limp(\psi\lor\eta)), \\
	(\phi\lor\psi\liff\eta)
	\defiff
	((\phi\lor\psi)\liff\eta), \\
	(\phi\liff\psi\lor\eta)
	\defiff
	(\phi\liff(\psi\lor\eta)), \\
	(\phi\limp\psi\liff\eta)
	\defiff
	((\phi\limp\psi)\liff\eta), \\
	(\phi\liff\psi\limp\eta)
	\defiff
	(\phi\liff(\psi\limp\eta)).
\end{gather}
\end{definition}

% \begin{definition}
% 设\(\phi,\psi\)都是合式公式.
% 定义：
% \((\neg(\phi\limp\psi))
% \defiff
% (\phi\lnimp\psi)\).
% \end{definition}

\begin{example}
\([(\neg p) \lor (q)]\)等同于\([\neg p \lor q]\).
    \begin{solution}
        否定的优先级高于析取，
        所以先计算\(\neg p\)，
        再与\(q\)作析取.
        删去所示括号不改变公式的结构.
    \end{solution}
\end{example}

\begin{example}
    \([p \limp q \land r \limp s]\)等同于\([p \limp ((q \land r) \limp s)]\).
    \begin{solution}
        先计算\(q\land r\)，
        再按同级运算从右向左的约定计算\((q\land r)\limp s\)，
        最后以\(p\)为前件计算整个条件式.
        左右两种结合方式一般不等值.
        例如在\((p,q,r,s)=(0,1,1,0)\)时，
        有
        \begin{align*}
            p\limp((q\land r)\limp s) &= 0\limp0=1, \\
            (p\limp(q\land r))\limp s &= 1\limp0=0.
        \end{align*}
    \end{solution}
\end{example}

\subsection{命题公式的类型}
\begin{definition}
%@see: 《离散数学》（邓辉文） P88 定义3-2
在任何指派下均取真的命题公式，
称为\DefineConcept{永真式}或\DefineConcept{重言式}（tautology）.
\end{definition}

\begin{definition}
%@see: 《离散数学》（邓辉文） P88 定义3-2
在任何指派下均取假的命题公式，
称为\DefineConcept{永假式}或\DefineConcept{矛盾式}（contradiction）.
\end{definition}

\begin{definition}
%@see: 《离散数学》（邓辉文） P88 定义3-2
至少有一种指派使其为真的命题公式，
称为\DefineConcept{可满足式}（satisfactable formula）.
\end{definition}

\begin{definition}
%@see: 《离散数学》（邓辉文） P88 定义3-2
至少有一种指派使其为真，同时至少有一种指派使其为假的命题公式，
称为\DefineConcept{中性式}或\DefineConcept{偶然式}（contingency）.
\end{definition}

\begin{example}
命题公式\(\neg p \lor q \limp r\)是中性式.
    \begin{proof}
        原式按优先级解释为\((\neg p\lor q)\limp r\).
        在\((p,q,r)=(1,0,0)\)时，
        前件为假，
        所以原式为真.
        在\((p,q,r)=(0,0,0)\)时，
        前件为真而后件为假，
        所以原式为假.
        故原式既有成真指派又有成假指派，
        是中性式.
    \end{proof}
\end{example}

\begin{example}
命题公式\(p \lor \neg p\)是永真式.
    \begin{proof}
        任取真值指派.
        若\(p\)为真，
        析取的第一项为真；若\(p\)为假，
        \(\neg p\)为真.
        两种情况下原式均为真，
        所以它是永真式.
    \end{proof}
\end{example}

\begin{example}
命题公式\(p \land \neg p\)是永假式.
    \begin{proof}
        在任意真值指派下，
        \(p\)与\(\neg p\)恰有一个为假，
        所以它们的合取总为假.
        故原式是永假式.
    \end{proof}
\end{example}

\begin{example}
%@see: 《离散数学》（邓辉文） P89 例3-9
证明：命题公式\((p \limp q) \liff (\neg p \lor q)\)是永真式.
\begin{proof}
列出\((p \limp q) \liff (\neg p \lor q)\)的真值表：\begin{center}
%@Mathematica: TableForm[BooleanTable[{p, q, Implies[p, q], Not[p], Or[Not[p], q], Equivalent[Implies[p, q], Or[Not[p], q]]}, {p, q}] // Boole, TableHeadings -> {None, {"p", "q", "p \[Implies] q", "\[Not]p", "\[Not]p\[Or]q", "(p\[Implies]q)\[Equivalent](\[Not]p\[Or]q)"}}]
	\begin{tblr}{*2c|*3c|c}
		\hline
		\(p\) & \(q\) & \(p \limp q\) & \(\neg p\) & \(\neg p \lor q\) & \((p \limp q) \liff (\neg p \lor q)\) \\
		\hline
		1 & 1 & 1 & 0 & 1 & 1 \\
		1 & 0 & 0 & 0 & 0 & 1 \\
		0 & 1 & 1 & 1 & 1 & 1 \\
		0 & 0 & 1 & 1 & 1 & 1 \\
		\hline
	\end{tblr}
\end{center}
由真值表可知，命题公式\((p \limp q) \liff (\neg p \lor q)\)是永真式.
\end{proof}
\end{example}
\begin{remark}
像这样，通过绘制真值表，得出一个命题公式的类型的方法，称为\emph{真值表法}.
真值表法是最常用的判断命题公式类型的方法，
但是话说回来，当命题变元较多时，真值表法是极为不方便的.
\end{remark}

\begin{example}
%@see: 《离散数学》（邓辉文） P89 例3-10
证明：命题公式\((p \land (p \limp q)) \limp q\)是永真式.
\begin{proof}
假设\(p \land (p \limp q)\)取真，
则\(p\)及\(p \limp q\)均取真，
进而\(q\)为真，
因此\((p \land (p \limp q)) \limp q\)永真.
\end{proof}
\end{example}
\begin{remark}
像这样，利用对取值进行分析，得出一个命题公式的类型的方法，称为\emph{取值法}.
\end{remark}

\begin{example}
证明：命题公式\(p \limp p\)是永真式.
% 同一律
\begin{proof}
列出\(p \limp p\)的真值表：\begin{center}
%@Mathematica: TableForm[BooleanTable[{p, Implies[p, p]}, p] // Boole, TableHeadings -> {None, {"p", "p \[Implies] p"}}]
	\begin{tblr}{c|c}
		\hline
		\(p\) & \(p \limp p\) \\
		\hline
		1 & 1 \\
		0 & 1 \\
		\hline
	\end{tblr}
\end{center}
由真值表可知，命题公式\(p \limp p\)是永真式.
\end{proof}
\end{example}

\begin{example}
%@see: 《离散数学》（邓辉文） P90 习题 4.(1)
证明：命题公式\(p \limp (q \limp p)\)是永真式.
% 肯定后件律
\begin{proof}
列出\(p \limp (q \limp p)\)的真值表：\begin{center}
%@Mathematica: TableForm[BooleanTable[{p, q, Implies[q, p], Implies[p, Implies[q, p]]}, {p, q}] // Boole, TableHeadings -> {None, {"p", "q", "q \[Implies] p", "p\[Implies](q \[Implies] p)"}}]
	\begin{tblr}{*2c|c|c}
		\hline
		\(p\) & \(q\) & \(q \limp p\) & \(p \limp (q \limp p)\) \\
		\hline
		1 & 1 & 1 & 1 \\
		1 & 0 & 1 & 1 \\
		0 & 1 & 0 & 1 \\
		0 & 0 & 1 & 1 \\
		\hline
	\end{tblr}
\end{center}
由真值表可知，命题公式\(p \limp (q \limp p)\)是永真式.
\end{proof}
\end{example}

\begin{example}
证明：命题公式\((p \limp (q \limp r)) \limp ((p \limp q) \limp (p \limp r))\)是永真式.
\begin{proof}
列出\((p \limp (q \limp r)) \limp ((p \limp q) \limp (p \limp r))\)的真值表：\begin{center}
%@Mathematica: TableForm[BooleanTable[{p, q, r, Implies[Implies[p, Implies[q, r]], Implies[Implies[p, q], Implies[p, r]]]}, {p, q, r}] // Boole, TableHeadings -> {None, {"p", "q", "r", "(p\[Implies](q\[Implies]r))\[Implies]((p\[Implies]q)\[Implies](p\\[Implies]r))"}}]
	\begin{tblr}{*3c|c}
		\hline
		\(p\) & \(q\) & \(r\) & \((p \limp (q \limp r)) \limp ((p \limp q) \limp (p \limp r))\) \\
		\hline
		1 & 1 & 1 & 1 \\
		1 & 1 & 0 & 1 \\
		1 & 0 & 1 & 1 \\
		1 & 0 & 0 & 1 \\
		0 & 1 & 1 & 1 \\
		0 & 1 & 0 & 1 \\
		0 & 0 & 1 & 1 \\
		0 & 0 & 0 & 1 \\
		\hline
	\end{tblr}
\end{center}
由真值表可知，命题公式\((p \limp (q \limp r)) \limp ((p \limp q) \limp (p \limp r))\)是永真式.
\end{proof}
\end{example}

\begin{example}
证明：命题公式\((\neg p \limp \neg q) \limp (q \limp p)\)是永真式.
\begin{proof}
列出\((\neg p \limp \neg q) \limp (q \limp p)\)的真值表：\begin{center}
%@Mathematica: TableForm[BooleanTable[{p, q, Implies[Implies[Not[p], Not[q]], Implies[q, p]]}, {p, q}] // Boole, TableHeadings -> {None, {"p", "q", "(\[Not]p\[Implies]\[Not]q)\[Implies](q\[Implies]p)"}}]
	\begin{tblr}{*2c|c}
		\hline
		\(p\) & \(q\) & \((\neg p \limp \neg q) \limp (q \limp p)\) \\
		\hline
		1 & 1 & 1 \\
		1 & 0 & 1 \\
		0 & 1 & 1 \\
		0 & 0 & 1 \\
		\hline
	\end{tblr}
\end{center}
由真值表可知，命题公式\((\neg p \limp \neg q) \limp (q \limp p)\)是永真式.
\end{proof}
\end{example}

\begin{example}
证明：\(p \limp q\)是\(p \land r \limp q \land r\)的充分不必要条件.
\begin{proof}
列出真值表：\begin{center}
%@Mathematica: TableForm[BooleanTable[{p, q, r, Implies[p, q], Implies[And[p, r], And[q, r]], Implies[Implies[p, q], Implies[And[p, r], And[q, r]]], Implies[Implies[And[p, r], And[q, r]], Implies[p, q]]}, {p, q, r}] // Boole, TableHeadings -> {None, {"p", "q", "r", "p\[Implies]q", "p\[And]r\[Implies]q\[And]r","(p\[Implies]q)\[Implies](p\[And]r\[Implies]q\[And]r)", "(p\[And]r\[Implies]q\[And]r)\[Implies](p\[Implies]q)"}}]
	\begin{tabular}{*3c|cc|cc}
		\hline
		\(p\) & \(q\) & \(r\)
			& \(p \limp q\)
			& \(p \land r \limp q \land r\)
			& \((p \limp q) \limp (p \land r \limp q \land r)\)
			& \((p \land r \limp q \land r) \limp (p \limp q)\) \\
		\hline
		1 & 1 & 1 & 1 & 1 & 1 & 1 \\
		1 & 1 & 0 & 1 & 1 & 1 & 1 \\
		1 & 0 & 1 & 0 & 0 & 1 & 1 \\
		1 & 0 & 0 & 0 & 1 & 1 & 0 \\
		0 & 1 & 1 & 1 & 1 & 1 & 1 \\
		0 & 1 & 0 & 1 & 1 & 1 & 1 \\
		0 & 0 & 1 & 1 & 1 & 1 & 1 \\
		0 & 0 & 0 & 1 & 1 & 1 & 1 \\
		\hline
	\end{tabular}
\end{center}
立即可知\(p \limp q\)是\(p \land r \limp q \land r\)的充分不必要条件.
\end{proof}
\end{example}

\subsection{永真式代入定理}
下面介绍\DefineConcept{永真式代入定理}（Rule of Substitution，简称RS）.
\begin{theorem}
%@see: 《离散数学》（邓辉文） P89 定理3-1
设命题公式\(A(\AutoTuple{p}{n})\)是永真式，
则分别用命题公式\(\AutoTuple{B}{n}\)代换命题变元\(\AutoTuple{p}{n}\)，
所得到的命题公式\(A(\AutoTuple{B}{n})\)也是永真式.
    \begin{proof}
        记\(v(F)\)为公式\(F\)在指派\(v\)下的真值.
        任取代入后公式的指派\(v\)，
        定义原公式变元上的指派\(w\)，
        使\(w(p_i)=v(B_i)\).
        按\(A\)的构造归纳可得
        \begin{align*}
            v\bigl(A(B_1,\dotsc,B_n)\bigr) &= w\bigl(A(p_1,\dotsc,p_n)\bigr).
        \end{align*}
        当\(A=p_i\)时，
        等式由\(w\)的定义成立；当\(A\)是常量时，
        两侧取同一常量.
        若等式对直接子公式成立，
        则对其真值施加同一个否定或二元联结词的真值函数后仍相等，
        故归纳成立.
        由于\(A\)永真，
        右侧为\(1\)，
        左侧也为\(1\).
        指派\(v\)任意，
        所以代入后的公式永真.
        这里采用同时代入，
        不要求各\(B_i\)之间不共享命题变元.
    \end{proof}
\end{theorem}
% \DefineConcept{代入实例}（substitution instance）

\section{逻辑等值的命题公式}
\subsection{逻辑等值的定义}
\begin{definition}
%@see: 《离散数学》（邓辉文） P90 定义3-3
给定两个命题公式\(A\)和\(B\)，
如果在任何真值指派下，\(A\)和\(B\)的真值都相同，
则称“命题公式\(A\)和\(B\) \DefineConcept{逻辑等值}（logically equal）”
或“命题公式\(A\)和\(B\) \DefineConcept{逻辑等价}（logically equivalent）”，
记为\(A = B\).
反之，如果在某个真值指派下，\(A\)和\(B\)的真值不同，
则记\(A \neq B\).
\end{definition}
\begin{remark}
注意不要混淆“逻辑等价”与“等价联结词”这两个概念.
\end{remark}

\begin{theorem}
%@see: 《离散数学》（邓辉文） P91 定理3-2
设\(A\)和\(B\)是命题公式，
则\(A = B\)的充分必要条件是\(A \liff B\)是永真式.
    \begin{proof}
        等价联结词的真值为\(1\)，
        当且仅当两个操作数的真值相同.
        因此对于任意指派\(v\)，
        有
        \begin{align*}
            v(A\liff B)=1 &\iff v(A)=v(B).
        \end{align*}
        对所有\(v\)要求上述条件成立，
        左侧正是\(A\liff B\)永真，
        右侧正是\(A=B\).
    \end{proof}
\end{theorem}

\begin{example}
%@see: 《离散数学》（邓辉文） P91 例3-11
证明：对于任意命题公式\(p,q\)，
有\begin{equation}\label{equation:数理逻辑.蕴含式化为析取式}
	p \limp q = \neg p \lor q.
\end{equation}
\begin{proof}
列出真值表：\begin{center}
	\begin{tblr}{*5c}
		\hline
		\(p\) & \(q\) & \(p \limp q\) & \(\neg p\) & \(\neg p \lor q\) \\
		\hline
		0 & 0 & 1 & 1 & 1 \\
		0 & 1 & 1 & 1 & 1 \\
		1 & 0 & 0 & 0 & 0 \\
		1 & 1 & 1 & 0 & 1 \\
		\hline
	\end{tblr}
\end{center}
显然有\(p \limp q = \neg p \lor q\).
\end{proof}
\end{example}

\begin{example}
证明：对于任意命题公式\(p,q,r\)，有\begin{equation}
	(p \limp r) \land (q \limp r) = (p \lor q) \limp r.
\end{equation}
%@Mathematica: TableForm[BooleanTable[{p, q, r, And[Implies[p, r], Implies[q, r]], Implies[Or[p, q], r]}, {p, q, r}] // Boole, TableHeadings -> {None, {"p", "q", "r", "(p\[Implies]r)\[And](q\[Implies]r)", "(p\[Or]q)\[Implies]r"}}]
\begin{proof}
列出真值表：\begin{center}
	\begin{tblr}{*5c}
		\hline
		\(p\) & \(q\) & \(r\)
		& \((p \limp r) \land (q \limp r)\)
		& \((p \lor q) \limp r\) \\ \hline
		1 & 1 & 1 & 1 & 1 \\
		1 & 1 & 0 & 0 & 0 \\
		1 & 0 & 1 & 1 & 1 \\
		1 & 0 & 0 & 0 & 0 \\
		0 & 1 & 1 & 1 & 1 \\
		0 & 1 & 0 & 0 & 0 \\
		0 & 0 & 1 & 1 & 1 \\
		0 & 0 & 0 & 1 & 1 \\
		\hline
	\end{tblr}
\end{center}
显然有\((p \limp r) \land (q \limp r) = (p \lor q) \limp r\).
\end{proof}
\end{example}

\begin{example}
证明：对于任意命题公式\(p,q,r\)，有\begin{equation}
	(p \limp q) \land (p \limp r) = p \limp (q \land r).
\end{equation}
%@Mathematica: TableForm[BooleanTable[{p, q, r, And[Implies[p, q], Implies[p, r]], Implies[p, And[q, r]]}, {p, q, r}] // Boole, TableHeadings -> {None, {"p", "q", "r", "(p\[Implies]q)\[And](p\[Implies]r)", "p\[Implies](q\[And]r)"}}]
\begin{proof}
列出真值表：\begin{center}
	\begin{tblr}{*5c}
		\hline
		\(p\) & \(q\) & \(r\)
		& \((p \limp q) \land (p \limp r)\)
		& \(p \limp (q \land r)\) \\ \hline
		1 & 1 & 1 & 1 & 1 \\
		1 & 1 & 0 & 0 & 0 \\
		1 & 0 & 1 & 0 & 0 \\
		1 & 0 & 0 & 0 & 0 \\
		0 & 1 & 1 & 1 & 1 \\
		0 & 1 & 0 & 1 & 1 \\
		0 & 0 & 1 & 1 & 1 \\
		0 & 0 & 0 & 1 & 1 \\
		\hline
	\end{tblr}
\end{center}
显然有\((p \limp q) \land (p \limp r) = p \limp (q \land r)\).
\end{proof}
\end{example}

%@Mathematica: TableForm[BooleanTable[{p, q, r, And[p, And[q, r]], And[And[p, q], r]}, {p, q, r}] // Boole, TableHeadings -> {None, {"p", "q", "r", "p\[And](q\[And]r)", "(p\[And]q)\[And]r"}}]
%@Mathematica: TableForm[BooleanTable[{p, q, r, Or[p, Or[q, r]], Or[Or[p, q], r]}, {p, q, r}] // Boole, TableHeadings -> {None, {"p", "q", "r", "p\[Or](q\[Or]r)", "(p\[Or]q)\[Or]r"}}]
%@Mathematica: TableForm[BooleanTable[{p, q, r, Xor[p, Xor[q, r]], Xor[Xor[p, q], r]}, {p, q, r}] // Boole, TableHeadings -> {None, {"p", "q", "r", "p\[Xor](q\[Xor]r)", "(p\[Xor]q)\[Xor]r"}}]

\begin{example}
    设\(p,q,r\)是互异的命题变元.
    证明下列两个公式不逻辑等值：\begin{equation*}
	p \limp (q \limp r)
	\neq
	(p \limp q) \limp r.
\end{equation*}
\begin{proof}
列出真值表：\begin{center}
%@Mathematica: TableForm[BooleanTable[{p, q, r, Implies[p, Implies[q, r]], Implies[Implies[p, q], r]}, {p, q, r}] // Boole, TableHeadings -> {None, {"p", "q", "r", "p\[Equivalent](q\[Equivalent]r)", "(p\[Equivalent]q)\[Equivalent]r"}}]
	\begin{tblr}{*3c|cc}
		\hline
		\(p\) & \(q\) & \(r\)
		& \(p \limp (q \limp r)\)
		& \((p \limp q) \limp r\) \\ \hline
		1 & 1 & 1 & 1 & 1 \\
		1 & 1 & 0 & 0 & 0 \\
		1 & 0 & 1 & 1 & 1 \\
		1 & 0 & 0 & 1 & 1 \\
		0 & 1 & 1 & 1 & 1 \\
		0 & 1 & 0 & 1 & 0 \\
		0 & 0 & 1 & 1 & 1 \\
		0 & 0 & 0 & 1 & 0 \\
		\hline
	\end{tblr}
\end{center}
显然有\(p \limp (q \limp r) \neq (p \limp q) \limp r\).
\end{proof}
\end{example}

\begin{example}
证明：对于任意命题公式\(p,q,r\)，有\begin{equation*}
	p \liff (q \liff r)
	= (p \liff q) \liff r.
\end{equation*}
\begin{proof}
列出真值表：\begin{center}
%@Mathematica: TableForm[BooleanTable[{p, q, r, Equivalent[p, Equivalent[q, r]], Equivalent[Equivalent[p, q], r]}, {p, q, r}] // Boole, TableHeadings -> {None, {"p", "q", "r", "p\[Equivalent](q\[Equivalent]r)", "(p\[Equivalent]q)\[Equivalent]r"}}]
	\begin{tblr}{*3c|cc}
		\hline
		\(p\) & \(q\) & \(r\)
		& \(p \liff (q \liff r)\)
		& \((p \liff q) \liff r\) \\ \hline
		1 & 1 & 1 & 1 & 1 \\
		1 & 1 & 0 & 0 & 0 \\
		1 & 0 & 1 & 0 & 0 \\
		1 & 0 & 0 & 1 & 1 \\
		0 & 1 & 1 & 0 & 0 \\
		0 & 1 & 0 & 1 & 1 \\
		0 & 0 & 1 & 1 & 1 \\
		0 & 0 & 0 & 0 & 0 \\
		\hline
	\end{tblr}
\end{center}
显然有\(p \liff (q \liff r) = (p \liff q) \liff r\).
\end{proof}
\end{example}

\begin{example}
    设\(p,q,r\)是互异的命题变元.
    证明下列两个公式不逻辑等值：\begin{equation*}
	p \lnand (q \lnand r)
	\neq
	(p \lnand q) \lnand r.
\end{equation*}
\begin{proof}
列出真值表：\begin{center}
%@Mathematica: TableForm[BooleanTable[{p, q, r, Nand[p, Nand[q, r]], Nand[Nand[p, q], r]}, {p, q, r}] // Boole, TableHeadings -> {None, {"p", "q", "r", "p\[Nand](q\[Nand]r)", "(p\[Nand]q)\[Nand]r"}}]
	\begin{tblr}{*3c|cc}
		\hline
		\(p\) & \(q\) & \(r\)
		& \(p \lnand (q \lnand r)\)
		& \((p \lnand q) \lnand r\) \\ \hline
		1 & 1 & 1 & 1 & 1 \\
		1 & 1 & 0 & 0 & 1 \\
		1 & 0 & 1 & 0 & 0 \\
		1 & 0 & 0 & 0 & 1 \\
		0 & 1 & 1 & 1 & 0 \\
		0 & 1 & 0 & 1 & 1 \\
		0 & 0 & 1 & 1 & 0 \\
		0 & 0 & 0 & 1 & 1 \\
		\hline
	\end{tblr}
\end{center}
显然有\(p \lnand (q \lnand r) \neq (p \lnand q) \lnand r\).
\end{proof}
\end{example}

\begin{example}
    设\(p,q,r\)是互异的命题变元.
    证明下列两个公式不逻辑等值：\begin{equation*}
	p \lnor (q \lnor r)
	\neq
	(p \lnor q) \lnor r.
\end{equation*}
\begin{proof}
列出真值表：\begin{center}
%@Mathematica: TableForm[BooleanTable[{p, q, r, Nor[p, Nor[q, r]], Nor[Nor[p, q], r]}, {p, q, r}] // Boole, TableHeadings -> {None, {"p", "q", "r", "p\[Nor](q\[Nor]r)", "(p\[Nor]q)\[Nor]r"}}]
	\begin{tblr}{*3c|cc}
		\hline
		\(p\) & \(q\) & \(r\)
		& \(p \lnor (q \lnor r)\)
		& \((p \lnor q) \lnor r\) \\ \hline
		1 & 1 & 1 & 0 & 0 \\
		1 & 1 & 0 & 0 & 1 \\
		1 & 0 & 1 & 0 & 0 \\
		1 & 0 & 0 & 0 & 1 \\
		0 & 1 & 1 & 1 & 0 \\
		0 & 1 & 0 & 1 & 1 \\
		0 & 0 & 1 & 1 & 0 \\
		0 & 0 & 0 & 0 & 0 \\
		\hline
	\end{tblr}
\end{center}
显然有\(p \lnor (q \lnor r) \neq (p \lnor q) \lnor r\).
\end{proof}
\end{example}

\begin{theorem}
%@see: 《离散数学》（邓辉文） P91 定理3-3
设\(A_1(\AutoTuple{p}{n}) = A_2(\AutoTuple{p}{n})\).
分别用命题公式\(\AutoTuple{B}{n}\)代换命题变元\(\AutoTuple{p}{n}\)，
所得到的两个命题公式等值，即\begin{equation*}
	A_1(\AutoTuple{B}{n}) = A_2(\AutoTuple{B}{n}).
\end{equation*}
    \begin{proof}
        由假设及前一定理，
        \(A_1(p_1,\dotsc,p_n)\liff A_2(p_1,\dotsc,p_n)\)永真.
        同时将\(p_i\)代入\(B_i\)，
        由永真式代入定理，
        \(A_1(B_1,\dotsc,B_n)\liff A_2(B_1,\dotsc,B_n)\)仍永真.
        再由逻辑等值与等价式永真的充分必要条件，
        即得结论.
    \end{proof}
\end{theorem}

\begin{theorem}
%@see: 《离散数学》（邓辉文） P91 定理3-4
对于任意命题公式\(A,B,C\)，有\begin{itemize}
	\item {\rm\bf 自反性}：\(A = A\).
	\item {\rm\bf 对称性}：若\(A = B\)，则\(B = A\).
	\item {\rm\bf 传递性}：若\(A = B\)且\(B = C\)，则\(A = C\).
\end{itemize}
    \begin{proof}
        任取指派\(v\).
        因为\(v(A)=v(A)\)，
        所以\(A=A\).
        若\(A=B\)，
        则\(v(A)=v(B)\)，
        也就有\(v(B)=v(A)\)，
        故\(B=A\).
        若\(A=B\)且\(B=C\)，
        则\(v(A)=v(B)=v(C)\)，
        故\(A=C\).
        这些结论对任意指派成立，
        依次给出自反性、对称性和传递性.
    \end{proof}
\end{theorem}

\subsection{基本等值式}
\begin{theorem}\label{theorem:数理逻辑.基本等值式}
%@see: 《离散数学》（邓辉文） P91 定理3-5
对于任意命题公式\(A,B,C\)，
有\begin{itemize}
	\item {\rm\bf 对合律}：\begin{equation}
		% 双重否定律
		\neg\neg A = A.
	\end{equation}

	\item {\rm\bf 幂等律}（或称{\rm\bf 重叠律}）：\begin{gather}
		A \lor A = A, \\
		A \land A = A.
	\end{gather}

	\item {\rm\bf 交换律}：\begin{gather}
		A \lor B = B \lor A, \\
		A \land B = B \land A.
	\end{gather}

	\item {\rm\bf 结合律}：\begin{gather}
		(A \lor B) \lor C = A \lor (B \lor C), \\
		(A \land B) \land C = A \land (B \land C).
	\end{gather}

	\item {\rm\bf 吸收律}：\begin{gather}
		A \lor (A \land B) = A, \\
		A \land (A \lor B) = A.
	\end{gather}

	\item {\rm\bf 分配律}：\begin{gather}
		A \lor (B \land C) = (A \lor B) \land (A \lor C), \\
		A \land (B \lor C) = (A \land B) \lor (A \land C).
	\end{gather}

	\item {\rm\bf 互补律}：\begin{gather}
		% 排中律
		A \lor \neg A = 1, \\
		% 矛盾律
		A \land \neg A = 0.
	\end{gather}

	\item {\rm\bf 德摩根律}：\begin{gather}
		\neg(A \lor B) = \neg A \land \neg B, \\
		\neg(A \land B) = \neg A \lor \neg B.
	\end{gather}

	\item {\rm\bf 同一律}：\begin{gather}
		A \lor 0 = 0 \lor A = A, \\
		A \land 1 = 1 \land A = A.
	\end{gather}

	\item {\rm\bf 0-1 律}：\begin{gather}
		A \lor 1 = 1 \lor A = 1, \\
		A \land 0 = 0 \land A = 0.
	\end{gather}
\end{itemize}
    \begin{proof}
        固定任意指派，
        记\(A,B,C\)的真值分别为\(a,b,c\).
        以下逐项验证各条等值式.

        两次否定得到原真值，
        即\(1-(1-a)=a\)，
        所以对合律成立.
        \(a\lor a\)与\(a\land a\)均等于\(a\)，
        所以幂等律成立.
        析取为真当且仅当至少一个操作数为真，
        合取为真当且仅当两个操作数均为真.
        这些条件与操作数次序无关，
        所以交换律成立.
        两种加括号方式的三项析取均在至少一项为真时取真，
        三项合取均在三项全真时取真，
        所以结合律成立.

        当\(a=0\)时，
        \(a\lor(a\land b)=0\)；当\(a=1\)时，
        该式为\(1\).
        同样，
        \(a\land(a\lor b)\)在\(a=0\)时为\(0\)，
        在\(a=1\)时为\(1\)，
        所以两条吸收律成立.
        第一条分配律在\(a=1\)时两侧均为\(1\)，
        在\(a=0\)时两侧均为\(b\land c\).
        第二条分配律在\(a=0\)时两侧均为\(0\)，
        在\(a=1\)时两侧均为\(b\lor c\)，
        所以两条分配律成立.

        \(a\)与\(1-a\)恰有一个为真，
        所以它们的析取为\(1\)，
        合取为\(0\)，
        得到互补律.
        \(\neg(A\lor B)\)为真当且仅当\(A,B\)均假，
        与\(\neg A\land\neg B\)为真的条件相同.
        \(\neg(A\land B)\)为真当且仅当\(A,B\)至少一个为假，
        与\(\neg A\lor\neg B\)为真的条件相同，
        得到德摩根律.

        与\(0\)析取或与\(1\)合取不改变\(a\)，
        得到同一律；与\(1\)析取总得\(1\)，
        与\(0\)合取总得\(0\)，
        得到\(0\)-\(1\)律.
        由于指派任意，
        所有等值式均成立.
    \end{proof}
\end{theorem}

\begin{theorem}
%@see: 《离散数学》（邓辉文） P92 定理3-6
对于任意命题公式\(A,B\)，有\begin{gather}
	A \lxor B
	= \neg(A \liff B),
		\label{equation:数理逻辑.联结词等值式1} \\
	A \limp B
	= \neg A \lor B,
		\label{equation:数理逻辑.联结词等值式2} \\
	A \liff B
	= (A \limp B) \land (B \limp A),
		\label{equation:数理逻辑.联结词等值式3} \\
	A \lnand B
	= \neg(A \land B),
		\label{equation:数理逻辑.联结词等值式4} \\
	A \lnor B
	= \neg(A \lor B),
		\label{equation:数理逻辑.联结词等值式5} \\
	A \lnimp B
	= \neg(A \limp B).
		\label{equation:数理逻辑.联结词等值式6}
\end{gather}
    \begin{proof}
        任取指派.
        异或为真恰在两个操作数真值不同时，
        等价为真恰在二者真值相同时，
        所以第一式成立.
        条件式与\(\neg A\lor B\)均仅在\(A=1,B=0\)时为假，
        所以第二式成立.
        两个相反方向的条件式同时为真，
        恰好排除了两个操作数取值不同的两种情况，
        所以第三式成立.
        与非、或非和条件否定的真值表分别由合取、析取和条件的真值表逐项取反得到，
        故其余三式成立.
    \end{proof}
\end{theorem}

\begin{theorem}
%@see: 《离散数学》（邓辉文） P94 习题3.4 5.
%@see: 《离散数学》（邓辉文） P94 习题3.4 6.
对于任意命题公式\(A,B,C\)，有\begin{gather}
	A \lxor B = B \lxor A, \\
	(A \lxor B) \lxor C = A \lxor (B \lxor C), \\
	A \limp B = \neg B \limp \neg A, \\
	\neg(A \liff B) = A \liff \neg B, \\
	\neg A = A \lnand A, \\
	\neg A = A \lnor A, \\
	A \land B = (A \lnand B) \lnand (A \lnand B), \\
	A \land B = (A \lnor A) \lnor (B \lnor B), \\
	A \lor B = (A \lnand A) \lnand (B \lnand B), \\
	A \lor B = (A \lnor B) \lnor (A \lnor B), \\
	A \liff B = B \liff A, \\
	(A \liff B) \liff C = A \liff (B \liff C), \\
	A \liff B = \neg(A \lxor B), \\
	A \liff B = (A \land B) \lor (\neg A \land \neg B).
\end{gather}
    \begin{proof}
        固定任意指派.
        异或仅判断两个真值是否不同，
        所以异或可交换.
        三项异或不论如何加括号，
        都在三个操作数中有奇数个为真时取真，
        所以异或可结合.
        等价仅判断两个真值是否相同，
        所以它也可交换.
        若\(C\)为真，
        两种三项等价式都归结为\(A\liff B\)；若\(C\)为假，
        两种三项等价式都归结为\(A\lxor B\)，
        所以等价可结合.
        将一个操作数取反恰好交换“相同”和“不同”，
        这同时证明\(\neg(A\liff B)=A\liff\neg B\)及等价与异或互为否定.
        两个操作数真值相同，
        恰指它们同时真或同时假，
        所以最后一式成立.
        由已证的基本等值式，
        还得到
        \begin{align*}
            A\limp B &= \neg A\lor B=B\lor\neg A=\neg B\limp\neg A, \\
            A\lnand A &= \neg(A\land A)=\neg A, \\
            A\lnor A &= \neg(A\lor A)=\neg A, \\
            (A\lnand B)\lnand(A\lnand B) &= \neg\neg(A\land B)=A\land B, \\
            (A\lnor A)\lnor(B\lnor B) &= \neg(\neg A\lor\neg B)=A\land B, \\
            (A\lnand A)\lnand(B\lnand B) &= \neg(\neg A\land\neg B)=A\lor B, \\
            (A\lnor B)\lnor(A\lnor B) &= \neg\neg(A\lor B)=A\lor B.
        \end{align*}
        这就逐项证明了全部等值式.
    \end{proof}
\end{theorem}

\subsection{等值演算法}
基本等值式有很多用途，
可以用于化简命题公式（即将命题公式化为一个与其等值的满足指定条件的含联结词最少的命题公式）、
判断命题公式的类型、证明等值式、计算命题公式的范式、命题逻辑推理等.

在使用等值式时，下列\DefineConcept{等值置换定理}（Rule of Replacement，简称RR）是至关重要的.
\begin{theorem}
设\(C\)是命题公式\(A\)的子公式，且\(C = D\)，
则将\(A\)中的\(C\)部分或全部替换成\(D\)所得到的命题公式与\(A\)等值.
    \begin{proof}
        先只替换\(C\)的一处出现，
        并固定任意指派.
        由于\(C=D\)，
        该位置替换前后的真值相同.
        沿公式的构造从该位置逐层向外考察：若外层是否定，
        相同真值取反后仍相同；若外层是二元联结词，
        另一个操作数不变，
        同一真值函数在相同输入上给出相同输出.
        有限次向外考察到整个公式，
        即知替换前后真值相同.
        指派任意，
        所以两式逻辑等值.
        部分或全部替换是有限次单处替换，
        再由等值关系的传递性即得结论.
    \end{proof}
\end{theorem}

利用基本等值式和等值置换定理求解问题的方法称为\emph{等值演算法}.

\begin{example}
%@see: 《离散数学》（邓辉文） P93 例3-13(1)
设\(A,B,C\)是任意的命题公式，
化简命题公式\((A \limp (B \lor \neg C)) \land \neg A \land B\)，
将最后结果表示为只含\(\neg\)和\(\lor\)的命题公式.
\begin{solution}
直接计算得\begin{align*}
	(A \limp (B \lor \neg C)) \land \neg A \land B
	&= (\neg A \lor (B \lor \neg C)) \land \neg A \land B \\
	&= ((\neg A \lor (B \lor \neg C)) \land \neg A) \land B \\
	% 吸收律
	&= \neg A \land B
	= \neg(A \lor \neg B).
\end{align*}
\end{solution}
\end{example}

\subsection{对偶原理}
在\cref{theorem:数理逻辑.基本等值式} 中，
除了对合律以外，其他性质都是成对出现的，两者间有一定的联系.
下面我们给出命题公式的对偶式的定义.
\begin{definition}
%@see: 《离散数学》（邓辉文） P94 定义3-4
    设命题公式\(A\)只使用\(\neg,\land,\lor\)这三种联结词，
    但不限制它们出现的次数.
将\(A\)中\(\land\)换成\(\lor\)，\(\lor\)换成\(\land\)，\(1\)换成\(0\)，\(0\)换成\(1\)，
所得到的命题公式称为“\(A\)的\DefineConcept{对偶式}（dual formula）”，记为\(A^*\).
\end{definition}

\begin{example}
%@see: 《离散数学》（邓辉文） P94 例3-16(1)
设\(p,q\)是命题变元，写出命题公式\(\neg(p \land q) \land 1\)的对偶式.
\begin{solution}
命题公式\(\neg(p \land q) \land 1\)的对偶式为
\(\neg(p \lor q) \lor 0\).
\end{solution}
\end{example}

\begin{example}
%@see: 《离散数学》（邓辉文） P94 例3-16(2)
设\(p,q,r\)是命题变元，写出命题公式\(p \lor (q \land r)\)的对偶式.
\begin{solution}
命题公式\(p \lor (q \land r)\)的对偶式为
\(p \land (q \lor r)\).
\end{solution}
\end{example}

\begin{theorem}[对偶原理]
    %@see: 《离散数学》（邓辉文） P94 定理3-8（对偶原理）
    设\(A,B\)只使用\(\neg,\land,\lor\)及命题常量、命题变元.
    若\(A=B\)，
    则\(A^*=B^*\).
    \begin{proof}
        固定指派\(v\)，
        定义互补指派\(\bar v\)，
        使每个变元\(p\)满足\(\bar v(p)=1-v(p)\).
        按公式\(F\)的构造归纳证明
        \begin{align*}
            v(F^*) &= 1-\bar v(F).
        \end{align*}
        对变元，
        等式由互补指派的定义成立；对常量，
        对偶交换\(0,1\)，
        等式也成立.
        若\(F=\neg G\)，
        由归纳假设有
        \begin{align*}
            v(F^*) &= 1-v(G^*)=\bar v(G)=1-\bar v(\neg G).
        \end{align*}
        若\(F=G\land H\)，
        则\(F^*=G^*\lor H^*\)，
        归纳假设与德摩根律说明其真值为\(\bar v(G\land H)\)的否定.
        若\(F=G\lor H\)，
        同理由另一条德摩根律得到所需等式.
        因此归纳完成.
        由\(A=B\)可知\(\bar v(A)=\bar v(B)\)，
        从而
        \begin{align*}
            v(A^*) &= 1-\bar v(A)=1-\bar v(B)=v(B^*).
        \end{align*}
        指派\(v\)任意，
        所以\(A^*=B^*\).
    \end{proof}
\end{theorem}

对偶原理有利于记忆\cref{theorem:数理逻辑.基本等值式}，
这是因为除了对合律以外的其余每一对性质，都只需要记住其中一个即可.

\section{命题公式的范式}
%@see: https://www.geeksforgeeks.org/normal-and-principle-forms/
%@see: https://library.fiveable.me/key-terms/formal-logic-ii/conjunctive-normal-form
给定一个命题公式，根据它的真值表，显然可以方便地得出它在每一种指派下的真值.
但是随着命题变元个数\(N\)的增加，命题公式的真值指派就有\(2^N\)种，在实际计算中就变成不可行的了.

\begin{definition}
设\(A\)是命题公式.
若\begin{equation*}
	A = A_1 \lor A_2 \lor \dotsb \lor A_n
	\quad(n\geq1),
\end{equation*}
其中\(A_i\ (i=1,2,\dotsc,n)\)是由命题变元或其否定组成的合取式，
则称“\(A_1 \lor A_2 \lor \dotsb \lor A_n\)是
\(A\)的\DefineConcept{析取范式}（disjunctive normal form）”.
\end{definition}

\begin{definition}
设\(A\)是命题公式.
若\begin{equation*}
	A = A_1 \land A_2 \land \dotsb \land A_n
	\quad(n\geq1),
\end{equation*}
其中\(A_i\ (i=1,2,\dotsc,n)\)是由命题变元或其否定组成的析取式，
则称“\(A_1 \land A_2 \land \dotsb \land A_n\)是
\(A\)的\DefineConcept{合取范式}（conjunctive normal form）”.
\end{definition}

\begin{theorem}
    任意命题公式都存在析取范式与合取范式.
    \begin{proof}
        先利用联结词等值式消去\(\neg,\land,\lor\)以外的联结词，
        再利用德摩根律和对合律将否定移到命题变元前面.
        常量\(1,0\)分别用\(p\lor\neg p\)、\(p\land\neg p\)表示；
        若原式没有变元，
        允许引入这个不影响真值的辅助变元\(p\).
        于是只需处理由文字（命题变元或其否定）经合取、析取组成的公式.
        单个文字本身同时是析取范式和合取范式.
        两个析取范式的析取仍是析取范式，
        它们的合取则用分配律展开为有限多个合取项的析取.
        这按公式结构归纳给出析取范式的存在性.
        对合取范式，
        两个合取范式的合取仍是合取范式，
        它们的析取则用另一条分配律展开为有限多个析取项的合取，
        所以合取范式也存在.
        原公式有限，
        每次展开只涉及有限多个项，
        因而构造在有限步后结束.
    \end{proof}
\end{theorem}

计算析取范式、合取范式的步骤：\begin{enumerate}
	\item 使用等值式，将命题公式中的联结词归约为\(\neg,\land,\lor\)；
	\item 利用德摩根律，将\(\neg\)移到命题变元的前面；
	\item 根据分配律，得到命题公式的析取范式、合取范式.
\end{enumerate}

\begin{example}
设\(p,q,r\)都是命题变元.
求命题公式\(A = p \limp q \liff r\)的析取范式、合取范式.
\begin{solution}
由等值式得\begin{align*}
	A &= p \limp q \liff r \\
	&= (\neg p \lor q \limp r) \land (r \limp \neg p \lor q) \\
	&= (\neg(\neg p \lor q) \lor r) \land (\neg r \lor (\neg p \lor q)) \\
	&= (p \land \neg q \lor r) \land (\neg p \lor q \lor \neg r).
\end{align*}
于是\(A\)的析取范式为\begin{align*}
	A &= (p \land \neg q \lor r) \land (\neg p \lor q \lor \neg r) \\
	% 分配律
	&= (p \land \neg q) \land (\neg p \lor q \lor \neg r)
		\lor r \land (\neg p \lor q \lor \neg r) \\
	% 分配律
	&= (p \land \neg q) \land \neg p
		\lor (p \land \neg q) \land q
		\lor (p \land \neg q) \land \neg r
		\lor r \land \neg p
		\lor r \land q
		\lor r \land \neg r \\
	% 吸收律，消去\(p \land \neg p\)、\(q \land \neg q\)和\(r \land \neg r\)
	&= p \land \neg q \land \neg r
		\lor \neg p \land r
		\lor q \land r,
\end{align*}
\(A\)的合取范式为\begin{align*}
	A &= (p \land \neg q \lor r) \land (\neg p \lor q \lor \neg r) \\
	% 分配律
	&= (p \lor r) \land (\neg q \lor r) \land (\neg p \lor q \lor \neg r).
\end{align*}
\end{solution}
\end{example}

根据命题公式的析取范式、合取范式，可以分别得出该命题公式的成真指派、成假指派.
例如，已知\(A\)的析取范式是
\(A = p \land \neg q \land \neg r
\lor \neg p \land r
\lor q \land r\)，
若要\(A\)取\(1\)，
则\begin{equation*}
	p \land \neg q \land \neg r,
	\qquad
	\neg p \land r
	\qquad
	q \land r
\end{equation*}至少一个为\(1\)，
由\(p \land \neg q \land \neg r = 1\)得\((p,q,r) = (1,0,0)\)，
由\(\neg p \land r = 1\)得\((p,q,r) = (0,1,1),(0,0,1)\)，
由\(q \land r = 1\)得\((p,q,r) = (1,1,1),(0,1,1)\)，
于是，\(A\)的成真指派有\begin{equation*}
	(1,0,0),
	(0,1,1),
	(0,0,1),
	(1,1,1).
\end{equation*}
同样地，由\(A\)的合取范式
\(A = (p \lor r) \land (\neg q \lor r) \land (\neg p \lor q \lor \neg r)\)可知，
\(A\)的成假指派有\begin{equation*}
	(1,0,1),
	(1,1,0),
	(0,1,0),
	(0,0,0).
\end{equation*}

一般来说，一个命题公式的析取范式、合取范式都不是唯一的.
例如，假设\(A = p\)，
那么\((p \land q) \lor (p \land \neg q)\)和\(p\)都是\(A\)的析取范式.
这种不唯一性给有些问题的讨论带来不便.
下面根据命题公式的所有命题变元，讨论给定命题公式的唯一的标准形式：
\DefineConcept{主析取范式}和\DefineConcept{主合取范式}.

\begin{definition}
对于给定的命题变元，若由命题变元或其否定组成的合取式满足\begin{itemize}
	\item 每个命题变元或其否定两者之一只出现一次，
	\item 按字典顺序或按下标从小到大顺序出现，
\end{itemize}
则称该合取式为“由所给命题变元产生的\DefineConcept{极小项}（minimal term）”.
\end{definition}

\begin{proposition}
    对固定的有限非空变元表\(p_1,\dotsc,p_n\)，
    每个极小项恰有一种指派使其取\(1\).
    \begin{proof}
        合取为真当且仅当每个文字都为真.
        若第\(i\)个文字是\(p_i\)，
        必须取\(p_i=1\)；若是\(\neg p_i\)，
        必须取\(p_i=0\).
        每个变元恰出现一次，
        所以这些条件唯一确定全部取值，
        且按此取值确实使每个文字为真.
    \end{proof}
\end{proposition}

根据这个结论，我们可以对极小项编码：
极小项用\(m_i\)表示，
其中下标\(i\)是由成真指派得到的二进制数或对应的十进制数.

\begin{table}[ht]
	\centering
	\begin{tblr}{c|c|c}
		\hline
		极小项 & 成真指派 & 极小项的符号表示\(m_i\) \\
		\hline
		\(p \land q\) & 11 & \(m_{11} = m_3\) \\
		\(p \land \neg q\) & 10 & \(m_{10} = m_2\) \\
		\(\neg p \land q\) & 01 & \(m_{01} = m_1\) \\
		\(\neg p \land \neg q\) & 00 & \(m_{00} = m_0\) \\
		\hline
	\end{tblr}
	\caption{由2个命题变元\(p,q\)产生的极小项及其成真指派、符号表示}
\end{table}

\begin{table}[ht]
	\centering
	\begin{tblr}{c|c|c}
		\hline
		极小项 & 成真指派 & 极小项的符号表示\(m_i\) \\
		\hline
		\(p \land q \land r\) & 111 & \(m_{111} = m_7\) \\
		\(p \land q \land \neg r\) & 110 & \(m_{110} = m_6\) \\
		\(p \land \neg q \land r\) & 101 & \(m_{101} = m_5\) \\
		\(p \land \neg q \land \neg r\) & 100 & \(m_{100} = m_4\) \\
		\(\neg p \land q \land r\) & 011 & \(m_{011} = m_3\) \\
		\(\neg p \land q \land \neg r\) & 010 & \(m_{010} = m_2\) \\
		\(\neg p \land \neg q \land r\) & 001 & \(m_{001} = m_1\) \\
		\(\neg p \land \neg q \land \neg r\) & 000 & \(m_{000} = m_0\) \\
		\hline
	\end{tblr}
	\caption{由3个命题变元\(p,q,r\)产生的极小项及其成真指派、符号表示}
\end{table}

\begin{definition}
对于命题公式\(A\)，
若由\(A\)中所有命题变元产生的若干个极小项的析取等值于\(A\)，
则把该析取式称为“\(A\)的\DefineConcept{主析取范式}（major disjunctive form）”.
\end{definition}
\begin{remark}
    本节的主范式相对于一个固定的有限非空变元表定义.
    讨论唯一性时约定各极小项或极大项不重复，
    并按编号递增排列；实际书写可以改变各项次序.
    “唯一”是在这种书写约定下而言；否则只能说忽略项的顺序和重复后唯一.
    对不含变元的常量公式可先固定一个辅助变元.

    全部\(2^n\)个极小项的析取是永真式\(1\)，
    所以永真式也有主析取范式.
    本节约定主范式至少含一项，
    因此永假式没有主析取范式，
    永真式没有主合取范式.
    若另行允许空析取为\(0\)、空合取为\(1\)，
    则这两个例外也可纳入，
    但不可混用两种约定.
\end{remark}

%@see: 《离散数学》（邓辉文） P99
利用\emph{等值演算法}，求命题公式\(A\)的主析取范式的计算步骤：\begin{enumerate}
	\item 求出\(A\)的析取范式；
	\item 利用分配律，补充析取范式缺少的命题变元.
\end{enumerate}

%@see: 《离散数学》（邓辉文） P100
利用\emph{真值表法}，求命题公式\(A\)的主析取范式的计算步骤：\begin{enumerate}
	\item 写出命题公式\(A\)的真值表；
	\item 写出\(A\)的成真指派对应的极小项，使得该极小项在该指派下也为\(1\)；
	\item \(A\)等值于所有这样写出的极小项的析取.
\end{enumerate}

\begin{definition}
对于给定的命题变元，若由命题变元或其否定组成的析取式满足\begin{itemize}
	\item 每个命题变元或其否定两者之一只出现一次，
	\item 按字典顺序或按下标从小到大顺序出现，
\end{itemize}
则称该析取式为“由所给命题变元产生的\DefineConcept{极大项}（maximal term）”.
\end{definition}

\begin{proposition}
    对固定的有限非空变元表\(p_1,\dotsc,p_n\)，
    每个极大项恰有一种指派使其取\(0\).
    \begin{proof}
        析取为假当且仅当每个文字都为假.
        若第\(i\)个文字是\(p_i\)，
        必须取\(p_i=0\)；若是\(\neg p_i\)，
        必须取\(p_i=1\).
        每个变元恰出现一次，
        所以这些条件唯一确定全部取值，
        且按此取值确实使每个文字为假.
    \end{proof}
\end{proposition}

根据这个结论，我们可以对极大项编码：
极大项用\(M_i\)表示，
其中下标\(i\)是由成假指派得到的二进制数或对应的十进制数.

\begin{table}[ht]
	\centering
	\begin{tblr}{c|c|c}
		\hline
		极大项 & 成假指派 & 极大项的符号表示\(M_i\) \\
		\hline
		\(p \lor q\) & 00 & \(M_{00} = M_0\) \\
		\(p \lor \neg q\) & 01 & \(M_{01} = M_1\) \\
		\(\neg p \lor q\) & 10 & \(M_{10} = M_2\) \\
		\(\neg p \lor \neg q\) & 11 & \(M_{11} = M_3\) \\
		\hline
	\end{tblr}
	\caption{由2个命题变元\(p,q\)产生的极大项及其成假指派、符号表示}
\end{table}

\begin{table}[ht]
	\centering
	\begin{tblr}{c|c|c}
		\hline
		极大项 & 成假指派 & 极大项的符号表示\(M_i\) \\
		\hline
		\(p \lor q \lor r\) & 000 & \(M_{000} = M_0\) \\
		\(p \lor q \lor \neg r\) & 001 & \(M_{001} = M_1\) \\
		\(p \lor \neg q \lor r\) & 010 & \(M_{010} = M_2\) \\
		\(p \lor \neg q \lor \neg r\) & 011 & \(M_{011} = M_3\) \\
		\(\neg p \lor q \lor r\) & 100 & \(M_{100} = M_4\) \\
		\(\neg p \lor q \lor \neg r\) & 101 & \(M_{101} = M_5\) \\
		\(\neg p \lor \neg q \lor r\) & 110 & \(M_{110} = M_6\) \\
		\(\neg p \lor \neg q \lor \neg r\) & 111 & \(M_{111} = M_7\) \\
		\hline
	\end{tblr}
	\caption{由3个命题变元\(p,q,r\)产生的极大项及其成假指派、符号表示}
\end{table}

\begin{definition}
对于命题公式\(A\)，
若由\(A\)中所有命题变元产生的若干个极大项的合取等值于\(A\)，
则把该合取式称为“\(A\)的\DefineConcept{主合取范式}（major conjunctive form）”.
\end{definition}

%@see: 《离散数学》（邓辉文） P102
利用\emph{等值演算法}，求命题公式\(A\)的主合取范式的计算步骤：\begin{enumerate}
	\item 求出\(A\)的合取范式；
	\item 利用分配律，补充合取范式缺少的命题变元.
\end{enumerate}

%@see: 《离散数学》（邓辉文） P102
利用\emph{真值表法}，求命题公式\(A\)的主合取范式的计算步骤：\begin{enumerate}
	\item 写出命题公式\(A\)的真值表；
	\item 写出\(A\)的成假指派对应的极大项，使得该极大项在该指派下也为\(0\)；
	\item \(A\)等值于所有这样写出的极大项的合取.
\end{enumerate}

\begin{example}
设命题公式\(A\)的真值表如下所示，求\(A\)的主析取范式与主合取范式.
\begin{center}
	\begin{tblr}{*4c|*4c}
		\hline
		\(p\) & \(q\) & \(r\) & \(A\) & \(p\) & \(q\) & \(r\) & \(A\) \\
		\hline
		1 & 1 & 1 & 1 & 0 & 1 & 1 & 0 \\
		1 & 1 & 0 & 0 & 0 & 1 & 0 & 0 \\
		1 & 0 & 1 & 0 & 0 & 0 & 1 & 0 \\
		1 & 0 & 0 & 1 & 0 & 0 & 0 & 1 \\
		\hline
	\end{tblr}
\end{center}
\begin{solution}
\(A\)等值于\(A\)的成真指派对应的极小项的析取：\begin{equation*}
	A = (p \land q \land r)
		\lor (p \land \neg q \land \neg r)
		\lor (\neg p \land \neg q \land \neg r).
\end{equation*}

\(A\)等值于\(A\)的成假指派对应的极大项的合取：\begin{equation*}
	A = (\neg p \lor \neg q \lor r)
		\land (\neg p \lor q \lor \neg r)
		\land (p \lor \neg q \lor \neg r)
		\land (p \lor \neg q \lor r)
		\land (p \lor q \lor \neg r).
\end{equation*}
\end{solution}
\end{example}
\begin{remark}
在电路实现逻辑运算时，通常选用项数较少的分解式.
在上例中，主析取范式的极小项个数为3，主合取范式的极大项个数为5，
因此最好选用主析取范式.
\end{remark}

\begin{theorem}
%@see: 《离散数学》（邓辉文） P102 定理3-9
任意非永假命题公式都存在唯一的主析取范式.
任意非永真命题公式都存在唯一的主合取范式.
    \begin{proof}
        固定包含原公式全部变元的非空有序变元表，
        采用前述无重复且按编号递增排列的约定.
        记\(T\)为\(A\)的全部成真指派的集合，
        \(F\)为全部成假指派的集合.
        对每个指派\(a\)，
        记\(m_a\)为仅在\(a\)下为真的极小项，
        \(M_a\)为仅在\(a\)下为假的极大项.

        若\(A\)非永假，
        则\(T\)非空，
        可以构造
        \begin{align*}
            D &= \bigvee_{a\in T}m_a.
        \end{align*}
        对于任何指派\(b\)，
        \(D\)为真当且仅当\(b\in T\)，
        这与\(A\)为真的条件相同，
        所以\(D=A\).
        任何与\(A\)等值的极小项析取都必须含有每个\(a\in T\)对应的\(m_a\)，
        否则在\(a\)下为假；也不能含\(a\notin T\)对应的项，
        否则在\(a\)下为真.
        故所含项恰由\(T\)确定，
        按约定排列后唯一.

        若\(A\)非永真，
        则\(F\)非空，
        同样构造
        \begin{align*}
            C &= \bigwedge_{a\in F}M_a.
        \end{align*}
        \(C\)恰在\(F\)中的指派下为假，
        故\(C=A\).
        遗漏\(F\)中的某项会使相应指派错误地取真，
        加入\(F\)外的某项会使相应指派错误地取假，
        所以项集由\(F\)唯一确定.
        这证明主合取范式的存在性和唯一性.
    \end{proof}
\end{theorem}

\section{联结词汇集的功能完备性}
本节用联结词集\(S\)表示公式时，
只允许从命题变元出发有限次使用\(S\)中的联结词，
不额外免费提供常量\(0,1\).
目标公式可以含常量，
但表示式必须构造出相应的恒定真值，
必要时允许引入不影响结果的辅助变元.
这个约定不可省略：若免费允许常量\(0\)，
则\(\neg p=p\limp0\)，
\(\{\land,\limp\}\)便不再是不完备的例子.
\begin{definition}
%@see: 《离散数学》（邓辉文） P106 定义3-11
设\(S\)是由联结词组成的汇集.
若任意命题公式都可由\(S\)中的联结词等值地表示出来，
则称“\(S\)是一个\DefineConcept{功能完备联结词集}（complete group of connectives）”.
\end{definition}

\begin{theorem}\label{theorem:数理逻辑.功能完备联结词集1}
%@see: 《离散数学》（邓辉文） P106 定理3-10
\(\{\neg,\land,\lor\}\)是功能完备联结词集.
    \begin{proof}
        本章各联结词均已表示为只含\(\neg,\land,\lor\)的等值式.
        按公式的构造逐层进行这些替换，
        由等值置换定理可得等值的表示式.
        常量\(1,0\)分别用\(p\lor\neg p\)、\(p\land\neg p\)表示.
        任意其他有限元真值联结词也能按真值表写成主析取范式，
        恒定真值的情况使用上述两个常量表示式.
        因此任意命题公式均可如此表示，
        该联结词集功能完备.
    \end{proof}
\end{theorem}
\begin{corollary}
\(\{\lnor\}\)是功能完备联结词集.
    \begin{proof}
        由已证的等值式，
        有
        \begin{align*}
            \neg A &= A\lnor A, \\
            A\lor B &= (A\lnor B)\lnor(A\lnor B), \\
            A\land B &= (A\lnor A)\lnor(B\lnor B).
        \end{align*}
        所以\(\neg,\land,\lor\)均可用\(\lnor\)表示.
        将前一定理中的表示式逐层替换，
        即得\(\{\lnor\}\)功能完备.
    \end{proof}
\end{corollary}
\begin{corollary}
\(\{\lnand\}\)是功能完备联结词集.
\begin{proof}
因为\begin{align*}
	\neg p
	&= \neg(p \land p)
	= p \lnand p, \\
	p \land q
	&= \neg(\neg(p \land q))
	= \neg(p \lnand q)
	= (p \lnand q) \lnand (p \lnand q), \\
	p \lor q
	&= \neg(\neg p \land \neg q)
	= (\neg p) \lnand (\neg q)
	= (p \lnand p) \lnand (q \lnand q),
\end{align*}
也就是说\(\neg p\)、\(p \land q\)、\(p \lor q\)
可以由仅含\(\lnand\)的命题公式等值表示，
那么根据\cref{theorem:数理逻辑.功能完备联结词集1} 可知，
\(\{\lnand\}\)是功能完备联结词集.
\end{proof}
\end{corollary}
\begin{corollary}
\(\{\neg,\land\}\)是功能完备联结词集.
    \begin{proof}
        由德摩根律，
        有
        \begin{align*}
            A\lor B &= \neg(\neg A\land\neg B).
        \end{align*}
        故\(\neg,\land,\lor\)均能用\(\neg,\land\)表示，
        由\cref{theorem:数理逻辑.功能完备联结词集1}及等值置换定理得到结论.
    \end{proof}
\end{corollary}
\begin{corollary}\label{theorem:数理逻辑.功能完备联结词集5}
\(\{\neg,\lor\}\)是功能完备联结词集.
\begin{proof}
因为\begin{align*}
	p \land q
	&= \neg(\neg p \lor \neg q),
\end{align*}
也就是说\(p \land q\)可由仅含\(\{\neg,\lor\}\)的命题公式等值表示，
那么根据\cref{theorem:数理逻辑.功能完备联结词集1} 可知，
\(\{\neg,\lor\}\)是功能完备联结词集.
\end{proof}
\end{corollary}
\begin{corollary}
\(\{\neg,\limp\}\)是功能完备联结词集.
    \begin{proof}
        条件式的析取表示及对合律给出
        \begin{align*}
            A\lor B &= \neg A\limp B.
        \end{align*}
        故\(\neg,\lor\)均能用\(\neg,\limp\)表示，
        由\cref{theorem:数理逻辑.功能完备联结词集5}及等值置换定理得到结论.
    \end{proof}
\end{corollary}

\begin{example}
%@see: 《离散数学》（邓辉文） P106 例3-26
定义3元联结词\(f\)：\begin{center}
	\begin{tblr}{*4c|*4c}
		\hline
		\(p\) & \(q\) & \(r\) & \(f(p,q,r)\) & \(p\) & \(q\) & \(r\) & \(f(p,q,r)\) \\
		\hline
		1 & 1 & 1 & 0 & 0 & 1 & 1 & 0 \\
		1 & 1 & 0 & 1 & 0 & 1 & 0 & 0 \\
		1 & 0 & 1 & 1 & 0 & 0 & 1 & 1 \\
		1 & 0 & 0 & 1 & 0 & 0 & 0 & 1 \\
		\hline
	\end{tblr}
\end{center}
证明：\(\{f\}\)是功能完备联结词集.
\begin{proof}
由\(f\)的真值表可知\(f(p,q,r)\)的主合取范式为\begin{equation*}
	f(p,q,r)
	= (\neg p \lor \neg q \lor \neg r)
	\land (p \lor \neg q \lor \neg r)
	\land (p \lor \neg q \lor r).
\end{equation*}
于是有\begin{equation*}
	f(p,p,p) = \neg p,
	\qquad
	f(\neg p,\neg p,\neg q) = p \lor q.
\end{equation*}
因此由\cref{theorem:数理逻辑.功能完备联结词集5} 可知，
\(\{f\}\)是功能完备联结词集.
\end{proof}
\end{example}

\begin{example}
    %@see: 《离散数学》（邓辉文） P107 例3-28
    在不免费使用命题常量的约定下，
    证明\(\{\land,\limp\}\)不是功能完备联结词集.
    \begin{proof}
        对只含这两个联结词的公式作结构归纳，
        证明在所有变元均取\(1\)时公式总取\(1\).
        原子公式是命题变元，
        结论成立.
        若结论对\(B,C\)成立，
        则全\(1\)指派下有\(B\land C=1\land1=1\)，
        \(B\limp C=1\limp1=1\)，
        所以归纳成立.
        但\(\neg p\)在全\(1\)指派下为\(0\)，
        不可能与上述任一公式等值，
        因而该联结词集不完备.
    \end{proof}
\end{example}

\begin{definition}
%@see: 《离散数学》（邓辉文） P107 定义3-12
设\(S\)是功能完备联结词集，
如果\(S\)的任意非空真子集都不是功能完备联结词集，
则称“\(S\)是\DefineConcept{极小功能完备联结词集}”.
\end{definition}

\begin{theorem}
%@see: 《离散数学》（邓辉文） P107 定理3-11
下列联结词集是极小功能完备的：\begin{itemize}
	\item \(\{\lnor\}\)；
	\item \(\{\lnand\}\)；
	\item \(\{\neg,\land\}\)；
	\item \(\{\neg,\lor\}\)；
	\item \(\{\neg,\limp\}\).
\end{itemize}
    \begin{proof}
        这五个集合的功能完备性已逐一证明.
        \(\{\lnor\}\)和\(\{\lnand\}\)没有非空真子集，
        符合定义中的极小性条件；即使检验空集，
        空集也只能产生命题变元，
        不能表示否定.

        对于其余三个集合，
        只需检验各自的两个单元素子集.
        只用\(\neg\)得到的公式至多依赖一个变元，
        不能表示同时实质依赖\(p,q\)的\(p\land q\)：固定\(p=1\)改变\(q\)，
        或固定\(q=1\)改变\(p\)，
        均能改变其取值.
        只用\(\land\)、只用\(\lor\)或只用\(\limp\)的公式，
        在全\(1\)指派下均取\(1\).
        这一点由原子变元及\(1\land1=1\)、\(1\lor1=1\)、\(1\limp1=1\)作结构归纳即得.
        所以这三个单元素集均不能表示\(\neg p\).
        全部非空真子集都不完备，
        因此所列集合均极小功能完备.
    \end{proof}
\end{theorem}

\section{命题逻辑中的推理}
逻辑学的主要内容是研究推理.
推理是从一些前提推出结论的思维过程.
实际问题中的推理，需要对前提做深入分析，才能得出结论.
例如由两直线平行，得出同位角相等；
以及由一元二次方程的判别式大于0，得出方程有两个不相等的实数根.

数理逻辑主要是用数学的方法研究逻辑中的推理，
它关心的是推理形式的有效性问题.

\subsection{推理形式有效性的定义}
\begin{definition}
%@see: 《离散数学》（邓辉文） P108 定义3-13
设\(\AutoTuple{H}{n}\)和\(C\)是命题公式.
若由\(\AutoTuple{H}{n}\)全为真，可得出\(C\)必然真，
则称“由\(\AutoTuple{H}{n}\)得出\(C\)的\DefineConcept{推理形式是有效的}（valid argument form）”
或“\(\AutoTuple{H}{n}\) \DefineConcept{逻辑推出}（logically follows） \(C\)”
或“\(\AutoTuple{H}{n}\) \DefineConcept{逻辑蕴含}（logically implies） \(C\)”，
记为\begin{equation*}
	\AutoTuple{H}{n} \implies C
	\quad\text{或}\quad
	\AutoTuple{H}{n} \models C
	~\footnote{
		当命题公式\(C\)是永真式时，
		不论\(\AutoTuple{H}{n}\)是否全为真，
		\(C\)恒为真，
		于是我们可以省略\(\AutoTuple{H}{n}\)，
		把\(\AutoTuple{H}{n} \models C\)
		简记为\(\models C\).
	};
\end{equation*}
把\(\AutoTuple{H}{n}\)称为\DefineConcept{前提}（antecedent，premise，hypothesis），
把\(C\)称为\DefineConcept{结论}（conclusion），
把\(\implies\)称为\DefineConcept{推出符号}，
把\(\AutoTuple{H}{n} \implies C\)称为\DefineConcept{推理规则}.
\end{definition}
% \(\vDash\)常用于形式系统的语义推理
% \(\vdash\)常用于形式系统的语构推理

\begin{theorem}\label{theorem:数理逻辑.推理规则的充分必要条件}
%@see: 《离散数学》（邓辉文） P109 定理3-12
设\(\AutoTuple{H}{n}\)和\(C\)是命题公式，
则\(\AutoTuple{H}{n} \implies C\)的充分必要条件是
    \(H_1 \land \dotsb \land H_n \limp C\)是永真式.
\begin{proof}
命题公式由\(\AutoTuple{H}{n}\)全为真，
        可得出\(H_1 \land \dotsb \land H_n\)全为真.
于是，由\(\AutoTuple{H}{n}\)全为真可得出\(C\)必然真的充分必要条件是
        \(H_1 \land \dotsb \land H_n \limp C\)是永真式.
\end{proof}
\end{theorem}
\begin{remark}
由于成立\cref{theorem:数理逻辑.推理规则的充分必要条件}，
所以我们又把\(\AutoTuple{H}{n} \implies C\)
称为\DefineConcept{永真蕴含式}或\DefineConcept{逻辑蕴含式}，
并相应地把\(\implies\)称为\DefineConcept{永真蕴含符号}或\DefineConcept{逻辑蕴含符号}.
但是一定要注意：\(\implies\)与蕴含连接词\(\limp\)是不同的.
就两条命题公式而言，\(\implies\)是关系符号，\(A \implies B\)是逻辑蕴含式，
\(\limp\)是运算符号，\(A \limp B\)是命题公式.
\end{remark}

\begin{definition}
设\(A\)和\(B\)是两个命题公式.
如果\(A = B\)，
则记\(A \iff B\)或\(A \dualmodels B\).
\end{definition}

\begin{theorem}
%@see: 《离散数学》（邓辉文） P109 定理3-13
设\(A\)和\(B\)是命题公式，
则\(A \iff B\)的充分必要条件是\(A \implies B\)且\(B \implies A\).
\begin{proof}
利用\(A \liff B = (A \limp B) \land (B \limp A)\)即得.
\end{proof}
\end{theorem}

若将公式视为符号串，
永真蕴含是自反且传递的预序关系，
不满足关于符号串相等的反对称性.
例如\(p\)与\(p\land p\)互相蕴含，
却不是同一个符号串.
下述定理中的\(A=B\)表示逻辑等值.
只有在按逻辑等值识别公式以后，
永真蕴含才给出偏序关系.
\begin{theorem}
%@see: 《离散数学》（邓辉文） P109 定理3-14
设\(A,B,C\)都是命题公式，则\begin{itemize}
	\item {\rm\bf 自反性}：\(A \implies A\)；
	\item {\rm\bf 反对称性}：如果\(A \implies B\)且\(B \implies A\)，则\(A = B\)；
	\item {\rm\bf 传递性}：如果\(A \implies B\)且\(B \implies C\)，则\(A \implies C\).
\end{itemize}
    \begin{proof}
        任取指派.
        \(A\)为真时\(A\)当然为真，
        故\(A\implies A\).
        若\(A\implies B\)且\(B\implies A\)，
        则\(A\)为真当且仅当\(B\)为真，
        所以二者真值相同，
        即\(A=B\).
        若\(A\implies B\)且\(B\implies C\)，
        每个使\(A\)为真的指派先使\(B\)为真，
        再使\(C\)为真，
        故\(A\implies C\).
    \end{proof}
\end{theorem}

\begin{theorem}
%@see: 《离散数学》（邓辉文） P109 定理3-15
设\(A,B,C\)都是命题公式，则\begin{itemize}
	\item 如果\(A \implies C\)且\(B \implies C\)，则\(A \lor B \implies C\)；
	\item 如果\(C \implies A\)且\(C \implies B\)，则\(C \implies A \land B\).
\end{itemize}
    \begin{proof}
        任取使\(A\lor B\)为真的指派，
        至少\(A,B\)之一为真.
        若\(A\)为真，
        由\(A\implies C\)得\(C\)为真；若\(B\)为真，
        由\(B\implies C\)也得\(C\)为真.
        所以第一条成立.
        对于第二条，
        任取使\(C\)为真的指派，
        由两个假设分别得\(A\)为真、\(B\)为真，
        故\(A\land B\)为真.
    \end{proof}
\end{theorem}

\begin{example}
    解释永真蕴含关系中的上确界\(\sup\)、下确界\(\inf\)，
    并求两个逻辑等值类的上确界和下确界.
    % 针对原有符号疑问的独立补充，并非未经核实的教材定理3-16转录.
    \begin{solution}
        记\([A]\)为所有与\(A\)逻辑等值的公式组成的类，
        定义\([A]\preccurlyeq[B]\)当且仅当\(A\implies B\).
        替换成逻辑等值的代表不改变任何指派下的真值，
        所以此关系与代表的选择无关.
        前一定理给出自反性、反对称性和传递性，
        因而它是偏序关系.

        两个类的上界是同时位于它们之上的类，
        上确界是其中不大于其他任何上界的类.
        下界与下确界反向定义.
        因为\(A\implies A\lor B\)、\(B\implies A\lor B\)，
        \([A\lor B]\)是上界.
        若\([C]\)是任一上界，
        则\(A\implies C\)、\(B\implies C\)，
        由前一定理得\(A\lor B\implies C\)，
        所以它是最小上界.
        同理\(A\land B\)同时蕴含\(A,B\)，
        且任一同时蕴含\(A,B\)的公式也蕴含\(A\land B\)，
        所以\([A\land B]\)是最大下界.
        因此
        \begin{align*}
            \sup\{[A],[B]\} &= [A\lor B], \\
            \inf\{[A],[B]\} &= [A\land B].
        \end{align*}
        若有两个最小上界，
        则它们互相不大于对方，
        由反对称性相等；最大下界的唯一性同理.
    \end{solution}
\end{example}

\subsection{基本推理规则}
下面举例说明，证明推理形式有效性的4种方法.

\begin{example}
%@see: 《离散数学》（邓辉文） P110 例3-29
设\(A,B\)是命题公式.
证明：\(A \limp B,A \implies B\).
\begin{proof}
\begin{proof}[证法一]
真值表法.
\begin{center}
%@Mathematica: TableForm[BooleanTable[{p, q, Implies[p, q], And[Implies[p, q], p], Implies[And[Implies[p, q], p], q]}, {p, q}] // Boole, TableHeadings -> {None, {"p", "q", "p \[Implies] q", "(p \[Implies] q)\[And]p", "((p \[Implies] q)\[And]p)\[Implies]q"}}]
	\begin{tblr}{*5c}
		\hline
		\(p\) & \(q\) & \(p \limp q\) & \((p \limp q) \land p\) & \(((p \limp q) \land p) \limp q\) \\
		\hline
		1 & 1 & 1 & 1 & 1 \\
		1 & 0 & 0 & 0 & 1 \\
		0 & 1 & 1 & 0 & 1 \\
		0 & 0 & 1 & 0 & 1 \\
		\hline
	\end{tblr}
\end{center}
由真值表可见\(((p \limp q) \land p) \limp q\)是永真式，
那么由\cref{theorem:数理逻辑.推理规则的充分必要条件} 可知\(A \limp B,A \implies B\).
\end{proof}
\begin{proof}[证法二]
取值法.
            假设\((p \limp q) \land p\)取真，
则\(p \limp q\)及\(p\)均取真，
进而\(q\)为真，
因此\(((p \limp q) \land p) \limp q\)永真.
\end{proof}
\begin{proof}[证法三]
等值演算法.
\begin{align*}
	((p \limp q) \land p) \limp q
	&= ((\neg p \lor q) \land p) \limp q \\
	&= ((\neg p \land p) \lor (p \land q)) \limp q \\
	&= (p \land q) \limp q \\
	&= \neg(p \land q) \lor q \\
	&= (\neg p \lor \neg q) \lor q \\
	&= \neg p \lor (\neg q \lor q) \\
	&= \neg p \lor 1
	= 1.
	\qedhere
\end{align*}
\end{proof}
\begin{proof}[证法四]
主范式法.
\(((p \limp q) \land p) \limp q\)的主析取范式为\begin{equation*}
	(p \land q) \lor (p \land \neg q) \lor (\neg p \land q) \lor (\neg p \land \neg q).
\end{equation*}
它的主合取范式不存在.
由主范式可见\(((p \limp q) \land p) \limp q\)永真.
\end{proof}
\let\qed\relax
\end{proof}
\end{example}

\begin{table}[hbt]
%@see: 《离散数学》（邓辉文） P110 表3-20
	\centering
	\begin{tblr}{cc|c}
		\hline
		化简 & Simplification
			& \(A \land B \implies A\) \\
		附加 & Addition
			& \(A \implies A \lor B\) \\
		合取引入 & Conjunction
			& \(A,B \implies A \land B\) \\
		析取三段论 & Disjunctive Syllogism
			& \(A \lor B, \neg A \implies B\) \\
		假言推理 & Modus Ponens
			& \(A \limp B, A \implies B\) \\
		拒取式 & Modus Tollens
			& \(A \limp B, \neg B \implies \neg A\) \\
		假言三段式 & Hypothetical Syllogism
			& \(A \limp B, B \limp C \implies A \limp C\) \\
		二难推理 & Constructive Dilemma
			& \(A \lor B, A \limp C, B \limp C \implies C\) \\
		% Destructive Dilemma
		\hline
	\end{tblr}
	\caption{基本逻辑蕴含式I}
\end{table}
\begin{proof}[基本逻辑蕴含式的证明]
    固定任意使每条规则全部前提为真的指派.
    合取为真时每项均真，
    得到化简规则；\(A\)为真时\(A\lor B\)为真，
    得到附加规则；\(A,B\)均真时合取为真，
    得到合取引入规则.
    若\(A\lor B\)真而\(A\)假，
    则\(B\)只能为真，
    得到析取三段论.
    若\(A\limp B\)及\(A\)均真，
    则\(B\)不能为假，
    得到假言推理.
    若\(A\limp B\)真而\(B\)假，
    则\(A\)只能为假，
    得到拒取式.
    若\(A\limp B\)、\(B\limp C\)均真，
    则\(A\)真时依次有\(B,C\)真，
    \(A\)假时\(A\limp C\)也真，
    得到假言三段式.
    最后，
    若\(A\lor B\)、\(A\limp C\)、\(B\limp C\)均真，
    则根据\(A,B\)中为真的那一项得到\(C\)为真，
    得到二难推理.
\end{proof}

\subsection{命题逻辑的自然推理系统}
自然推理的构造法是判定推理形式有效性的又一种方法，
它主要是为今后进一步学习数理逻辑，尤其是为逻辑的公理化推理系统做准备的.
自然推理的基本思想是，确定一些推理规则，然后根据这些推理规则从前提出发，把结论推出来.

作为推理系统，原则上有以下4个部分：\begin{enumerate}
	\item 推理系统应该有初始符号，即一些在推理系统中允许出现的字符，
	包括命题变元（通常用小写拉丁字母表示）、
	五个联结词（即\(\neg,\land,\lor,\limp,\liff\)）
	和左右圆括号；
	\item 合式公式的定义；
	\item 公理的确立；
	\item 推理规则的确立.
\end{enumerate}

在自然推理系统中不存在公理（这一点与公理推理系统截然不同）.
它把所有与5个联结词\(\neg,\land,\lor,\limp,\liff\)有关的基本逻辑蕴含式都作为推理规则.
同时，一个基本等值式相当于两个基本逻辑蕴含式.
除了基本逻辑蕴含式和基本等值式以外，还有两个最基本的推理规则：\begin{itemize}
	\item {\rm\bf P规则}，即所给的前提在证明过程中随时可以引用；
	\item {\rm\bf T规则}，即已经推出的命题公式在以后得证明过程中可以随时引用.
\end{itemize}

由于在自然推理系统中，作为推理依据的，只有推理规则，
这似乎更符合人们日常思维的推理习惯，因此称为自然推理.

\begin{table}[hbt]
%@see: 《离散数学》（邓辉文） P111 表3-21
	\centering
	\begin{tblr}{c|c}
		\hline
		对合律 & \(\neg\neg A = A\) \\
		幂等律 & \(A \lor A = A, A \land A = A\) \\
		交换律 & \(A \lor B = B \lor A, A \land B = B \land A\) \\
		结合律 & \((A \lor B) \lor C = A \lor (B \lor C), (A \land B) \land C = A \land (B \land C)\) \\
		吸收律 & \(A \lor (A \land B) = A, A \land (A \lor B) = A\) \\
		分配律 & \(A \lor (B \land C) = (A \lor B) \land (A \lor C), A \land (B \lor C) = (A \land B) \lor (A \land C)\) \\
		互补律 & \(A \lor \neg A = 1, A \land \neg A = 0\) \\
		De Morgan 律 & \(\neg(A \lor B) = \neg A \land \neg B, \neg(A \land B) = \neg A \lor \neg B\) \\
		同一律 & \(A \lor 0 = 0 \lor A = A, A \land 1 = 1 \land A = A\) \\
		0-1 律 & \(A \lor 1 = 1 \lor A = 1, A \land 0 = 0 \land A = 0\) \\
		& \(A \limp B = \neg A \lor B\) \\
		& \(A \liff B = (A \limp B) \land (B \limp A)\) \\
		\hline
	\end{tblr}
	\caption{基本等值式E}
\end{table}

在进行自然推理时，采用构造性证明方法
（简称\DefineConcept{构造法}或\DefineConcept{演绎法}或\DefineConcept{形式证明法}）.

通过以下例子了解自然推理证明过程的书写格式.
\begin{example}
%@see: 《离散数学》（邓辉文） P112 例3-30
证明：\(p \limp (q \lor r), \neg s \limp \neg q, p \land \neg s \implies r\).
\begin{proof}
用构造法.
\begin{center}
	\begin{tblr}{cp{4cm}l}
		(1) & \(p \land \neg s\) & P \\
		(2) & \(p\) & T(1)I \\
		(3) & \(\neg s\) & T(1)I \\
		(4) & \(p \limp (q \lor r)\) & P \\
		(5) & \(q \lor r\) & T(2)(4)I \\
		(6) & \(\neg s \limp \neg q\) & P \\
		(7) & \(\neg q\) & T(3)(6)I \\
		(8) & \(r\) & T(5)(7)I \\
	\end{tblr}
\end{center}
\end{proof}
\end{example}
\begin{remark}
从上述证明过程可以看出，每一行由3部分组成：
第一部分是编号，说明这一行是证明的第几步；
第二部分仅给出一个命题公式；
第三部分是给出理由，交代该命题公式是怎么得来的.
\end{remark}

下面介绍两种间接的构造性证明方法.

首先介绍反证法.
要证明\(\AutoTuple{H}{n} \implies C\)，
将结论\(C\)否定得到\(\neg C\)，然后推出一个矛盾（例如\(S \land \neg S\)）即可.

这个证明方法的道理是：要证明\(\AutoTuple{H}{n} \implies C\)，
只要证\(H_1 \land \dotsb \land H_n \limp C\)永真，
即\(\neg(H_1 \land \dotsb \land H_n) \lor C\)永真，
或\(\neg(\neg(H_1 \land \dotsb \land H_n) \lor C)\)永假.
\begin{example}
%@see: 《离散数学》（邓辉文） P113 例3-33
证明：\(p \land \neg q \limp r, \neg r \land \neg q \implies \neg p\).
\begin{proof}
用反证法.
\begin{center}
	\begin{tblr}{cp{4cm}l}
		(1) & \(\neg(\neg p)\) & P(附加) \\
		(2) & \(p\) & T(1)E \\
		(3) & \(\neg r \land \neg q\) & P \\
		(4) & \(\neg q\) & T(3)I \\
		(5) & \(p \land \neg q\) & T(2)(4)I \\
		(6) & \(p \land \neg q \limp r\) & P \\
		(7) & \(r\) & T(5)(6)I \\
		(8) & \(\neg r\) & T(3)I \\
		(9) & \(r \land \neg r\) & T(7)(8)I \\
	\end{tblr}
\end{center}
\end{proof}
\end{example}

下面介绍条件证明规则.
对于形如\(\AutoTuple{H}{n} \implies A \limp C\)的推理，
只需要证明\(\AutoTuple{H}{n},A \implies C\)，
这就是采用\DefineConcept{条件证明规则}的证明方法，
它是将\(A \limp C\)的前件\(A\)和后件\(C\)分离的一种证明方法.

\begin{proof}[条件证明规则的依据]
    记\(H=H_1\land\dotsb\land H_n\).
    不把逻辑推出关系当作命题公式参与等值运算，
    而只对条件式作以下计算：
    \begin{align*}
        H\limp(A\limp C) &= \neg H\lor(\neg A\lor C) \\
        &= (\neg H\lor\neg A)\lor C \\
        &= \neg(H\land A)\lor C \\
        &= (H\land A)\limp C.
    \end{align*}
    所以左式永真当且仅当右式永真.
    由\cref{theorem:数理逻辑.推理规则的充分必要条件}，
    这正表示\(H_1,\dotsc,H_n\implies A\limp C\)
    当且仅当\(H_1,\dotsc,H_n,A\implies C\).
\end{proof}

\begin{example}
%@see: 《离散数学》（邓辉文） P114 例3-34
证明：\(p \limp (q \lor r), q \limp \neg p, s \limp \neg r \implies p \limp \neg s\).
\begin{proof}
用条件证明规则.
\begin{center}
	\begin{tblr}{cp{4cm}l}
		(1) & \(p \limp (q \lor r)\) & P \\
		(2) & \(p\) & P(附加) \\
		(3) & \(q \lor r\) & T(1)(2)I \\
		(4) & \(q \limp \neg p\) & P \\
		(5) & \(\neg q \lor \neg p\) & T(4)E \\
		(6) & \(\neg q\) & T(2)(5)I \\
		(7) & \(r\) & T(3)(6)I \\
		(8) & \(s \limp \neg r\) & P \\
		(9) & \(\neg s \lor \neg r\) & T(8)E \\
		(10) & \(\neg s\) & T(7)(9)I \\
		(11) & \(p \limp \neg s\) & CP \\
	\end{tblr}
\end{center}
\end{proof}
\end{example}

数理逻辑的主要研究任务是建立一个严密的逻辑推理系统 --- 公理推理系统 --- 来刻画人类的思维规律，
这个系统与上述自然推理系统是类似的，但它有更精简的初始符号、公式的形成规则、公理和推理规则.
已有的逻辑演算形式系统，例如\emph{命题演算的形式系统}（formal system of proposition calculus），
在理论上证明了其合理性、一致性、完备性等与语法、语义有关的重要结论，
可以得出人类思维的所有推理规则，
所提供的逻辑推理框架能保证在前提为真的条件下，总能得出正确的结论.

\chapter{谓词逻辑}
原子命题是命题逻辑研究的基本单位，
没有对原子命题的内部结构及其逻辑关系进行讨论.
在实际思维中，仅有命题逻辑工具是不够的.
例如著名的苏格拉底三段论：
大前提是“所有的人都是要死的”，
小前提是“苏格拉底是人”，
结论是“所以，苏格拉底是要死的”.
这个推理的有效性在命题逻辑中无法证明，
因为上述每个命题都是原子命题，
可以分别用\(p,q,r\)表示，
然而\(p,q \implies r\)在命题逻辑中是无效推理.

之所以出现这种推理本身是正确的，但无法证明其有效性的问题，
是因为没有对原子命题的内部形式结构及其逻辑关系进行讨论，
这正是\DefineConcept{谓词逻辑}（predicate logic）
首先要研究的内容，这些讨论涉及集合、映射、运算和关系.

本章讨论的谓词逻辑又称为一阶逻辑，
它不允许对谓词、函词进行量化.

\section{个体、谓词、量词和函词}
\subsection{个体}
在谓词逻辑中，命题的考虑对象称为\DefineConcept{个体}（individual）.
个体是独立存在的事物.
个体可以是具体的，也可以是抽象的.

表示特定的、具体的个体，称为\DefineConcept{个体常量}（constant）.

不确定的个体，称为\DefineConcept{个体变元}（variable）.

个体常量和个体变元通常用小写拉丁字母表示，
例如\(a,b,c,\dotsc,x,y,z\).

在讨论个体时，通常要制定个体讨论的范围，
称为\DefineConcept{个体域}（domain of individuals）
或\DefineConcept{论域}（domain of discourse，universe），
用\(D\)表示，
本章始终假定\(D\)非空.

个体域可以是有限汇集，也可以是无限汇集.
我们把人类能够反映、认知的一切事物的汇集，
称为\DefineConcept{全总个体域}，简称\DefineConcept{全域}.

\subsection{谓词}
我们把表示个体性质以及个体之间关系的词称为\DefineConcept{谓词}（predicate）.

表示一个个体性质的谓词称为1元谓词.
表示\(n\ (n\geq0)\)个个体之间关系的谓词称为\(n\)元谓词.
谓词通常用大写拉丁字母表示，例如\(P,Q,R\).
为了强调\(n\)元谓词的元数，
会像表示\(n\)元函数一样，
在大写拉丁字母后面添上圆括号，并在括号内添上\(n\)个个体变元，
例如\(P(\AutoTuple{x}{n})\).

对于\(n\)元谓词\(P(\AutoTuple{x}{n})\ (n\geq1)\)，
当个体变元\(\AutoTuple{x}{n}\)取定个体域中的元素后，就成为一个命题，
例如\begin{equation*}
	G(x,y):
	x>y,
\end{equation*}
它是关于命题的函数，称为\DefineConcept{命题函数}（propositional function）.
显然，命题函数不是命题.

\begin{remark}
谓词的选取与个体域有关.
例如，对于命题“所有人都是要死的”，
若在所有人组成的个体域中考虑，
只需一个谓词\begin{equation*}
	D(x):
	\text{$x$是要死的};
\end{equation*}
若在全域中考虑，则需要两个谓词\begin{equation*}
	P(x):
	\text{$x$是人},
	\qquad
	D(x):
	\text{$x$是要死的},
\end{equation*}
其中\(P(x)\)称为\DefineConcept{特性谓词}，
我们用这个谓词将“人”从全域中分离出来.
\end{remark}

\subsection{量词}
对于命题函数，例如\begin{equation*}
	P(x):
	\text{$x$是素数},
\end{equation*}
当个体域\(D\)为自然数集\(\mathbb{N}\)时，
对于\(x\)的每一个取值，就有一个命题.
这是使\(P(x)\)成为命题的一种方法.

使\(P(x)\)成为命题的另一种方法是，
量化个体变元\(x\)，
常用方法有两种：全称量化和存在量化.

把表示个体数量特征的词称为\DefineConcept{量词}（quantifier）.
常见的量词有\DefineConcept{全称量词}~\(\color{blue}\forall\)（universal quantifier）、
\DefineConcept{存在量词}~\(\color{blue}\exists\)（existential quantifier）、
\DefineConcept{唯一量词}~\(\color{blue}\exists!\).

我们可以把存在量词转化为全称量词，
也就是说\(
	(\exists x)
	[\phi(x)]
\)
相当于\(
	\neg(\forall x)
	[\neg\phi(x)]
\).

\begin{proof}[存在量词转换的证明]
    固定论域、谓词解释和其余自由变元的取值.
    存在一个个体使\(\phi(x)\)为真，
    当且仅当并非每个个体都使\(\neg\phi(x)\)为真.
    这恰是前述两个公式各自为真的条件，
    所以它们逻辑等值.
\end{proof}

\begin{example}
    在带等号的一阶语言中，
    如何消去唯一量词\(\exists!\)？
    \begin{solution}
        此处\(x=y\)表示两个个体相同，
        不是表示两个公式逻辑等值.
        取未在\(\phi\)中出现的新变元\(y\)，
        以\(y\)替换\(x\)的自由出现，
        得到\(\phi(y)\).
        所需的表示是
        \begin{align*}
            (\exists!x)[\phi(x)] &= (\exists x)[\phi(x)\land(\forall y)(\phi(y)\limp y=x)].
        \end{align*}
        若左侧为真，
        取唯一满足\(\phi\)的个体作为\(x\)的见证，
        每个满足\(\phi(y)\)的个体必与它相同，
        故右侧为真.
        若右侧为真，
        则其见证满足\(\phi\)，
        且所有满足\(\phi\)的个体都等于该见证，
        所以存在性和唯一性同时成立.

        不能用“\(\phi(x)\)与\(\phi(y)\)真值相同则\(x=y\)”代替唯一性条件.
        例如取\(D=\{0,1,2\}\)，
        令\(\phi\)仅对\(0\)成立.
        唯一性成立，
        但\(\phi(1),\phi(2)\)同为假而\(1\ne2\).
        若原语言不含等号，
        就不能直接使用\(x=y\)作为原子公式；上述表示用于加入个体同一关系的语言.
    \end{solution}
\end{example}

由于在本章中，仅对个体进行量化，不对谓词进行量化，
因此把本章介绍的谓词逻辑特别称为\DefineConcept{一阶谓词逻辑}.

在使用量词时，必须要指明量化的个体变元，
例如\(\forall x,\forall y,\exists x,\exists y\).
单独使用量词是没有意义的.
我们把量词后面跟着的个体变元称为\DefineConcept{指导变元}.

若将命题函数中的所有个体变元都进行了量化，则得到一个命题，否则不是命题.

量词是对个体变元进行量化，所给的个体域\(D\)至关重要.
\begin{example}
    令\(G(x,y)\)表示\(x>y\).
    比较“没有最小的个体”在自然数域与整数域中的真值.
    \begin{solution}
        正确的量化形式是
        \begin{align*}
            (\forall x)(\exists y)[G(x,y)].
        \end{align*}
        在自然数域中，
        取最小自然数作为\(x\)，
        就不存在更小的自然数\(y\)，
        所以该式为假.
        在整数域中，
        对每个\(x\)取\(y=x-1\)即可，
        所以该式为真.
        \((\exists x)(\exists y)G(x,y)\)仅表示存在一对大小不同的个体，
        在两个域中都可取\(x=2,y=1\)，
        并不表示没有最小个体.
    \end{solution}
\end{example}

我们把量词\(\forall\)或\(\exists\)的作用范围
称为它的\DefineConcept{作用域}或\DefineConcept{辖域}（scope），
其辖域内受这个量词约束的、与指导变元同名的出现称为约束出现，
相应变元称为\DefineConcept{约束变元}（bound variable）.
辖域内其他名字的变元不因此受到约束.

若量词后有括号，则括号里面的部分是它的辖域.

若令\begin{equation*}
	P(x):
	\text{$x$是人},
	\qquad
	D(x):
	\text{$x$是要死的},
\end{equation*}
则“所有的人都是要死的”可以表示为\begin{equation*}
	(\forall x)[P(x) \limp D(x)],
\end{equation*}
这时\(\forall x\)的辖域为\(P(x) \limp D(x)\)，
其中两次出现的\(x\)是约束变元.

不受任何量词约束的变元称为\DefineConcept{自由变元}（free variable）.
例如\begin{equation*}
	(\forall x)[G(x,y)]
\end{equation*}中的\(y\)，
它不受\(\forall x\)的约束，这时\(G(x,y)\)中的\(y\)是自由变元.

\subsection{函词}
个体变元的函数称为\DefineConcept{函词}（function）.

例如，令个体域\(D\)取所有人组成的集合，用\(f(x)\)表示“\(x\)的父亲”，
则\(f\)是从\(D\)到\(D\)的一元函词.

又例如，令\(D \defeq \mathbb{R}\)，再令\(f(x,y) \defeq x^2 + y^2\)，
则\(f\)是从\(D^2\)到\(D\)的二元函词.

\section{谓词公式及命题的符号化}
\subsection{谓词公式}
\begin{example}
    给出一般一阶语言中项与合式谓词公式的递归定义，
    并说明自由出现、约束出现和代入的限制.
    \begin{solution}
        固定个体变元、个体常量以及各元数的函词、谓词符号.
        个体变元、个体常量都是\DefineConcept{项}；
        若\(f\)是\(n\)元函词且\(t_1,\dotsc,t_n\)是项，
        则\(f(t_1,\dotsc,t_n)\)是项.
        项仅由上述规则有限次生成.

        若\(P\)是\(n\)元谓词且\(t_1,\dotsc,t_n\)是项，
        则\(P(t_1,\dotsc,t_n)\)是\DefineConcept{原子公式}.
        \(0\)元谓词对应命题变元；在带等号的语言中，
        \(t_1=t_2\)也是原子公式.
        每个原子公式都是\DefineConcept{合式谓词公式}.
        若\(\phi,\psi\)为合式公式，
        则下列各式也是合式公式：
        \begin{align*}
            (\neg\phi),\quad(\phi\land\psi),\quad(\phi\lor\psi),\quad
            (\phi\limp\psi),\quad(\phi\liff\psi).
        \end{align*}
        其他联结词可以按已给出的等值式视为缩写.
        若\(\phi\)为合式公式且\(x\)为任一个体变元，
        则下列两式也为合式公式：
        \begin{align*}
            (\forall x)[\phi],\qquad(\exists x)[\phi].
        \end{align*}
        合式公式仅由这些规则有限次生成.
        \(a\in b\)只是选取二元谓词\(\in\)所得的特例，
        不能替代一般一阶语言的原子公式定义.

        自由或约束是某次出现的性质，
        不是预先把变元划为两种互斥的名字类别.
        量词\((\forall x)\)或\((\exists x)\)约束其辖域内原来自由的\(x\)，
        不要求\(x\)在加量词之前已经是约束变元.
        同一个名字在同一公式中可以有自由出现和约束出现.

        以项\(t\)替换\(x\)时，
        只替换\(x\)的自由出现，
        并要求\(t\)中的变元不会因此被同名量词捕获.
        满足此条件称为\(t\)对\(x\)可自由代入.
        若有冲突，
        先将相应约束变元一致地改名为新变元.
        例如将\((\forall y)R(x,y)\)中的自由\(x\)替换为\(y\)，
        应先改名得到\((\forall z)R(y,z)\)，
        而不是\((\forall y)R(y,y)\).
        以下\(A(t)\)均表示这种无捕获的代入.
    \end{solution}
\end{example}

%@see: 《离散数学》（邓辉文） P122
不含自由变元的谓词公式，称为\DefineConcept{闭式}（closed forumla）；
反之，含有自由变元的谓词公式，称为\DefineConcept{开式}（open forumla）.

\subsection{命题的符号化}
在一阶谓词逻辑中，将命题符号化时，必须使用谓词，并按以下步骤进行：
第一步是找出所给命题中的所有个体常量，分别符号化；
第二步是确定给定个体域中应该选用的所有谓词，特别注意特性谓词的选取；
第三步是确定量词、函词、联结词.

% \subsection{逻辑公理，推理规则}
% 我们有如下五条逻辑公理：
% \begin{axiom}
% \((\forall x)[\phi \limp \psi(x)] \limp [\phi \limp (\forall x) \psi(x)]\)，
% 其中，我们量化的自由变元\(a\)没有出现在公式\(\phi\)中.
% %where the free variable \(a\) on which we are quantifying does not occur in \(\phi\).
% \end{axiom}
% \begin{axiom}
% \((\forall x) \phi(x) \limp \phi(a)\)，
% 其中，\(\phi(a)\)是通过用自由变元\(a\)代替\(\phi(x)\)中的约束变元\(x\)得到的公式.
% %where \(\phi(a)\) is the forumla
% %obtained by replacing each occurrence
% %of the bound variable x
% %in \(\phi(x)\) by the free variable \(a\).
% \end{axiom}
% 以及如下两条\DefineConcept{推理规则}（rules of inference）：
% \begin{axiom}
% 从\(\phi\)和\(\phi \limp \psi\)，可以推断\(\psi\).
% %From \(\phi\) and \(\phi \limp \psi\) to infer \(\psi\).
% \end{axiom}
% \begin{axiom}
% 从\(\phi\)，可以推断\((\forall x) \phi(x)\)，
% 其中，\(\phi(x)\)表示在合式公式\(\phi\)中，
% 用约束变元\(x\)代替某个自由变元从而得到的公式.
% %From \(\phi\) to infer \((\forall x)\) \phi(x)
% %where \(\phi(x)\) is obtained from \(\phi\)
% %by replacing each occurrence of some free variable by x.
% \end{axiom}

\section{谓词公式的解释及类型}
\subsection{谓词公式的解释}
谓词公式的取值，取决于对它进行的解释或赋值，有如对命题公式的指派.
但与命题公式不同的时，谓词公式的解释有无限多种，
每种\DefineConcept{解释}（interpretation）由下面5个部分组成：\begin{enumerate}
	\item 指定个体域；
	\item 为谓词公式中的命题变元指派其真值；
	\item 将谓词公式中的个体常量及其自由变元解释为指定个体域中的元素；
	\item 将谓词公式中的函词解释为个体域上的函数；
	\item 将谓词公式中的谓词解释为个体域上的谓词.
\end{enumerate}

当把谓词公式中的个体常量\(a\)解释为个体域\(D\)中的个体\(\alpha\)时，
记作\(\frac{a}{\alpha}\).

谓词公式在任何解释下都会取得一个真值.
假设个体域是\(D\)，
且\(D\)是有限的
\footnote{
	根据定义，谓词公式是\emph{有限的}符号序列，所以不可以无限合取或无限析取.
	%@see: https://plato.stanford.edu/entries/logic-infinitary/
}，
则有以下两个消去量词的逻辑等值式：\begin{gather}
	(\forall x)[P(x)]
	= \bigwedge_{d \in D} P(d), \\
	(\exists x)[P(x)]
	= \bigvee_{d \in D} P(d).
\end{gather}
\begin{proof}
    将有限非空论域列为\(D=\{d_1,\dotsc,d_m\}\).
    全称式为真当且仅当每个\(P(d_i)\)为真，
    这恰是所列合取为真的条件.
    存在式为真当且仅当至少一个\(P(d_i)\)为真，
    这恰是所列析取为真的条件.
    此处\(P(d_i)\)表示在该解释中将\(x\)赋值为个体\(d_i\)所得的真值.
    若原语言没有为每个个体命名的常量，
    这只是语义层面的简写，
    并非声称每个\(d_i\)自动成为原语言中的符号.
\end{proof}

\subsection{谓词公式的类型}
%@see: 《离散数学》（邓辉文） P127 定义4-1
在任何解释下均为真的谓词公式称为\DefineConcept{永真式}或\DefineConcept{有效式}（valid）.

\begin{example}\label{example:数理逻辑.全称实例化}
    %@see: 《离散数学》（邓辉文） P127 例4-9
    设项\(t\)对\(x\)在\(A\)中可自由代入.
    证明\((\forall x)A(x)\limp A(t)\)永真.
    \begin{proof}
        固定任意解释及自由变元赋值，
        记\(t\)表示的个体为\(d\).
        若前件为假，
        条件式为真.
        若前件为真，
        \(A\)在\(x\)取任意个体时为真，
        特别在\(x=d\)时为真.
        无捕获代入保证\(A(t)\)具有这个真值，
        所以后件为真.
        两种情况下条件式均真，
        因而它永真.
    \end{proof}
\end{example}

\begin{example}
    %@see: 《离散数学》（邓辉文） P127 例4-10
    证明\(((\forall x)A(x)\lor(\forall x)B(x))\limp(\forall x)[A(x)\lor B(x)]\)永真.
    \begin{proof}
        固定任意解释及其余自由变元赋值.
        若前件为假，
        条件式为真.
        若前件为真，
        至少有一个全称式为真.
        若所有个体都满足\(A\)，
        则每个个体都满足\(A\lor B\)；若所有个体都满足\(B\)，
        结论相同.
        因而后件为真，
        整个条件式永真.
    \end{proof}
\end{example}

\begin{example}
%@see: 《离散数学》（邓辉文） P127 例4-11
证明谓词公式\(
	(\exists x)
	(\forall y)
	[A(x,y)]
	\limp
	(\forall y)
	(\exists x)
	[A(x,y)]
\)永真.
    \begin{proof}
        固定任意解释及其余自由变元赋值.
        若前件为真，
        则存在同一个个体\(a\)，
        使每个\(b\)均满足\(A(a,b)\).
        对\(y\)的任意取值\(b\)，
        选\(x=a\)即可，
        所以后件为真.
        前件为假时条件式也为真，
        故原式永真.
    \end{proof}
\end{example}

\begin{example}
%@see: 《离散数学》（邓辉文） P128 习题4.3 7.
证明谓词公式\(
	(\exists x)
	[
		A(x) \lor B(x)
	]
	\limp
	(\exists x)
	[A(x)]
	\lor
	(\exists x)
	[B(x)]
\)永真.
    \begin{proof}
        固定任意解释及其余自由变元赋值.
        若前件为真，
        取见证\(a\)，
        则\(A(a)\)或\(B(a)\)为真.
        前者使\((\exists x)A(x)\)为真，
        后者使\((\exists x)B(x)\)为真，
        所以后件的析取为真.
        前件为假时条件式也为真，
        故原式永真.
    \end{proof}
\end{example}

\begin{example}
%@see: 《离散数学》（邓辉文） P128 习题4.3 8.(1)
证明谓词公式\(
	(\forall x)
	[A(x)]
	\limp
	(\exists x)
	[A(x)]
\)永真.
    \begin{proof}
        固定任意解释.
        论域非空，
        可取其中一个个体\(a\).
        若\((\forall x)A(x)\)为真，
        则\(A(a)\)为真，
        所以\((\exists x)A(x)\)为真.
        前件为假时条件式也为真.
        非空论域条件不可省略：若允许空域，
        前件为真而后件为假.
    \end{proof}
\end{example}

\begin{example}
%@see: 《离散数学》（邓辉文） P128 习题4.3 8.(2)
证明谓词公式\(
	(\forall x)
	(\forall y)
	[A(x,y)]
	\limp
	(\forall x)
	(\exists y)
	[A(x,y)]
\)永真.
    \begin{proof}
        固定任意解释，
        假定前件为真.
        从非空论域中取一个\(b\).
        对\(x\)的任意取值\(a\)，
        前件保证\(A(a,b)\)为真，
        故\((\exists y)A(a,y)\)为真.
        由于\(a\)任意，
        后件为真.
        前件为假时条件式也真，
        所以原式永真.
    \end{proof}
\end{example}

\begin{example}
%@see: 《离散数学》（邓辉文） P128 习题4.3 8.(3)
证明谓词公式\(
	(\forall x)
	(\exists y)
	[A(x,y)]
	\limp
	(\exists x)
	(\exists y)
	[A(x,y)]
\)永真.
    \begin{proof}
        固定任意解释，
        假定前件为真.
        从非空论域中取一个\(a\)，
        由前件知存在\(b\)使\(A(a,b)\)为真.
        这对\(a,b\)就是后件的见证，
        所以后件为真.
        前件为假时条件式也真，
        故原式永真.
    \end{proof}
\end{example}

%@see: 《离散数学》（邓辉文） P127 定义4-2
至少存在一种解释使其成真的谓词公式称为\DefineConcept{可满足式}，
否则称为\DefineConcept{不可满足式}或\DefineConcept{矛盾式}或\DefineConcept{永假式}.

既存在取\(1\)的解释，又存在取\(0\)的解释的谓词公式
称为\DefineConcept{中性式}或\DefineConcept{偶然式}.

\begin{theorem}[半判定与判定的区别]
    对可有效编码的一般经典一阶语言，
    有效闭式与不可满足闭式均可半判定：对属于相应类别的输入，
    过程在有限步后给出肯定答案，
    对其余输入不保证停机.
    这不等于存在一个对所有输入均停机的判定过程，
    也不等于每个具体中性公式都无法判明.
    \begin{proof}
        这里调用一阶逻辑的可靠性与完备性定理：存在有限证明可有效检验的演算，
        其可证闭式恰好是全部有效闭式.
        按有限字符串长度及编码次序枚举所有候选证明，
        逐个检验，
        发现结论为给定闭式\(A\)的证明时即停机并报告“有效”.
        可靠性保证不误报，
        完备性保证每个有效\(A\)最终会被发现，
        因而这是半判定过程.
        又因为\(A\)不可满足当且仅当\(\neg A\)有效，
        对\(\neg A\)使用该过程即可半判定不可满足性.

        一般的一阶有效性没有通用判定算法，
        这是这里引用的丘奇不可判定性定理，
        不是从上述证明枚举直接得出的结论\footnote{
        A. Church，
        \emph{A note on the Entscheidungsproblem}，
        Journal of Symbolic Logic 1 (1936)，
        40--41，
        DOI: \href{https://doi.org/10.2307/2269326}{10.2307/2269326}，
        并参见同年勘误.
        关于可靠性、完备性与不可判定性，
        参见\href{https://slc.openlogicproject.org/}{\emph{Sets, Logic, Computation}}.
        }.
        若不可满足性可判定，
        则用\(\neg A\)是否不可满足便能判定\(A\)是否有效，
        同样与该定理矛盾.
        这里讨论不限制一般公式输入的判定问题，
        并不否认某些受限片段可判定.

        另一方面，
        取闭式\(P(c)\).
        在单元素论域中把\(c\)解释为该元素，
        分别令\(P\)成立和不成立，
        就得到一真一假两个解释，
        因而直接判明\(P(c)\)是中性式.
        这反驳了“每个中性谓词公式都不能在有限步内判明”的说法.
    \end{proof}
\end{theorem}

\section{逻辑等值的谓词公式}
\subsection{谓词公式等值的定义}
与两个命题公式等值完全类似，我们可以给出两个谓词公式等值的定义.
\begin{definition}
%@see: 《离散数学》（邓辉文） P129 定义4-3
设\(A,B\)是谓词公式.
若\(A\)和\(B\)在任何解释下的取值都相同，
则称“\(A\)和\(B\)是逻辑等值的”，
记为\(A=B\).
\end{definition}

\begin{proposition}
    \(A=B\)当且仅当\(A\liff B\)永真.
    \begin{proof}
        对每个解释及自由变元赋值，
        \(A\liff B\)为真当且仅当\(A,B\)真值相同.
        对全部解释及赋值要求此条件成立，
        就分别得到永真与逻辑等值的定义，
        故二者等价.
    \end{proof}
\end{proposition}

根据命题逻辑中的等值式，容易得到一些谓词逻辑中的等值式.

\subsection{基本等值式}
\subsubsection{量词转换}
%@see: 《离散数学》（邓辉文） P129
%@see: 《离散数学》（邓辉文） P131 习题4.4 1.
设谓词公式\(A\)含自由变元\(x\)，
则\begin{gather}
	\neg(\forall x)[A(x)]
	= (\exists x)[\neg A(x)],
	\label{equation:数理逻辑.量词转换1} \\
	\neg(\exists x)[A(x)]
	= (\forall x)[\neg A(x)].
	\label{equation:数理逻辑.量词转换2}
\end{gather}
\begin{proof}
    固定论域、谓词解释和其余自由变元的赋值.
    并非每个个体使\(A\)为真，
    当且仅当至少一个个体使\(A\)为假，
    所以第一式成立.
    不存在使\(A\)为真的个体，
    当且仅当每个个体使\(A\)为假，
    所以第二式成立.
    这些条件在所有解释及赋值下均等价，
    故为逻辑等值式.
\end{proof}


\subsubsection{量词辖域的收缩与扩张}
%@see: 《离散数学》（邓辉文） P129
%@see: 《离散数学》（邓辉文） P131 习题4.4 2.
设谓词公式\(A\)含自由变元\(x\)，
谓词公式\(B\)中不含自由变元\(x\)，
则\begin{gather}
	(\forall x)[A(x) \land B]
	= (\forall x)[A(x)] \land B, \\
	(\forall x)[A(x) \lor B]
	= (\forall x)[A(x)] \lor B, \\
	(\exists x)[A(x) \land B]
	= (\exists x)[A(x)] \land B, \\
	(\exists x)[A(x) \lor B]
	= (\exists x)[A(x)] \lor B.
\end{gather}
\begin{proof}
    固定解释及除\(x\)外的自由变元赋值.
    \(B\)没有自由的\(x\)，
    所以其真值不随\(x\)改变.
    若\(B\)真，
    则\(A(x)\land B\)与\(A(x)\)真值相同，
    而\(A(x)\lor B\)对每个个体均真.
    此时两条合取等式直接成立，
    两条析取等式两侧均真，
    其中存在式为真用到了非空论域.
    若\(B\)假，
    则\(A(x)\lor B\)与\(A(x)\)真值相同，
    而\(A(x)\land B\)对每个个体均假.
    此时两条析取等式直接成立，
    两条合取等式两侧均假，
    其中全称式为假也用到了非空论域.
    两种情况穷尽\(B\)的取值，
    故四式均成立.
\end{proof}


\begin{example}
%@see: 《离散数学》（邓辉文） P130 例4-13(1)
    设\(B\)中没有自由的\(x\).
证明：\begin{equation}
	(\forall x)[A(x) \limp B]
	= (\exists x)[A(x)] \limp B.
\end{equation}
\begin{proof}
显然\begin{align*}
	(\forall x)[A(x) \limp B]
	&= (\forall x)[\neg A(x) \lor B] \\
	&= (\forall x)[\neg A(x)] \lor B \\
	&= \neg(\exists x)[A(x)] \lor B \\
	&= (\exists x)[A(x)] \limp B.
	\qedhere
\end{align*}
\end{proof}
\end{example}

\begin{example}
%@see: 《离散数学》（邓辉文） P130 例4-13(2)
    设\(B\)中没有自由的\(x\).
证明：\begin{equation}
	(\forall x)[B \limp A(x)]
	= B \limp (\forall x)[A(x)].
\end{equation}
\begin{proof}
显然\begin{align*}
	(\forall x)[B \limp A(x)]
	&= (\forall x)[\neg B \lor A(x)] \\
	&= \neg B \lor (\forall x)[A(x)] \\
	&= B \limp (\forall x)[A(x)].
	\qedhere
\end{align*}
\end{proof}
\end{example}

\begin{example}
%@see: 《离散数学》（邓辉文） P130 例4-13(3)
%@see: 《离散数学》（邓辉文） P131 习题4.4 3.(1)
    设\(B\)中没有自由的\(x\).
证明：\begin{equation}
	(\exists x)[A(x) \limp B]
	= (\forall x)[A(x)] \limp B.
\end{equation}
\begin{proof}
显然\begin{align*}
	(\exists x)[A(x) \limp B]
	&= (\exists x)[\neg A(x) \lor B] \\
	&= (\exists x)[\neg A(x)] \lor B \\
	&= \neg(\forall x)[A(x)] \lor B \\
	&= (\forall x)[A(x)] \limp B.
	\qedhere
\end{align*}
\end{proof}
\end{example}

\begin{example}
%@see: 《离散数学》（邓辉文） P130 例4-13(4)
%@see: 《离散数学》（邓辉文） P131 习题4.4 3.(2)
    设\(B\)中没有自由的\(x\).
证明：\begin{equation}
	(\exists x)[B \limp A(x)]
	= B \limp (\exists x)[A(x)].
\end{equation}
\begin{proof}
显然\begin{align*}
	(\exists x)[B \limp A(x)]
	&= (\exists x)[\neg B \lor A(x)] \\
	&= \neg B \lor (\exists x)[A(x)] \\
	&= B \limp (\exists x)[A(x)].
	\qedhere
\end{align*}
\end{proof}
\end{example}

\subsubsection{量词分配律}
%@see: 《离散数学》（邓辉文） P130
%@see: 《离散数学》（邓辉文） P131 习题4.4 4.
设谓词公式\(A,B\)都含自由变元\(x\)，
则\begin{gather}
	(\forall x)[A(x) \land B(x)]
	= (\forall x)[A(x)] \land (\forall x)[B(x)], \\
	(\exists x)[A(x) \lor B(x)]
	= (\exists x)[A(x)] \lor (\exists x)[B(x)].
\end{gather}
\begin{proof}
    固定任意解释及其余自由变元赋值.
    每个个体同时满足\(A,B\)，
    当且仅当每个个体满足\(A\)并且每个个体满足\(B\)，
    故第一式两个方向均成立.
    若有个体满足\(A\lor B\)，
    则这个见证至少满足\(A,B\)之一，
    所以至少一个存在式为真.
    反之，
    任一存在式的见证都满足\(A\lor B\)，
    所以第二式两个方向也成立.
\end{proof}

\begin{remark}
量词\(\forall\)对\(\land\)可分配，
但是\((\forall x)[A(x) \lor B(x)]
\neq (\forall x)[A(x)] \lor (\forall x)[B(x)]\).
例如，令个体域\(D=\mathbb{Z}\)，
令\(A(x): \text{$x$是偶数}\)，
令\(B(x): \text{$x$是奇数}\)，
则\((\forall x)[A(x) \lor B(x)]\)在上述解释下取真，
而\((\forall x)[A(x)] \lor (\forall x)[B(x)]\)在上述解释下取假.

同样地，
量词\(\exists\)对\(\lor\)可分配，
但是\((\exists x)[A(x) \land B(x)]
\neq (\exists x)[A(x)] \land (\exists x)[B(x)]\).
\end{remark}
\begin{example}
    用有限论域说明全称量词一般不对析取分配，
    存在量词一般不对合取分配.
    \begin{solution}
        取\(D=\{0,1\}\)，
        令\(A\)仅对\(0\)为真，
        \(B\)仅对\(1\)为真.
        每个个体至少满足一个谓词，
        但两个谓词都不是处处真，
        所以\((\forall x)(A(x)\lor B(x))\)为真，
        \((\forall x)A(x)\lor(\forall x)B(x)\)为假.
        同时两个谓词各有见证却没有共同见证，
        所以\((\exists x)(A(x)\land B(x))\)为假，
        \((\exists x)A(x)\land(\exists x)B(x)\)为真.
        这分别给出了反例.
    \end{solution}
\end{example}

\begin{example}
%@see: 《离散数学》（邓辉文） P131 习题4.4 5.
证明：\begin{equation}
	(\exists x)[A(x) \limp B(x)]
	= (\forall x)[A(x)] \limp (\exists x)[B(x)].
\end{equation}
\begin{proof}
显然\begin{align*}
	(\exists x)[A(x) \limp B(x)]
	&= (\exists x)[\neg A(x) \lor B(x)] \\
	&= (\exists x)[\neg A(x)] \lor (\exists x)[B(x)] \\
	&= \neg(\forall x)[A(x)] \lor (\exists x)[B(x)] \\
	&= (\forall x)[A(x)] \limp (\exists x)[B(x)].
	\qedhere
\end{align*}
\end{proof}
\end{example}

\subsubsection{双重量词}
%@see: 《离散数学》（邓辉文） P130
设谓词公式\(A\)含自由变元\(x\)和\(y\)，
则\begin{gather}
	(\forall x)(\forall y)[A(x,y)]
	= (\forall y)(\forall x)[A(x,y)], \\
	(\exists x)(\exists y)[A(x,y)]
	= (\exists y)(\exists x)[A(x,y)].
\end{gather}
\begin{proof}
    固定任意解释及其余自由变元赋值.
    两个全称式均表示\(A\)对论域中的每一对个体\((a,b)\)为真，
    枚举第一、第二坐标的先后不影响这一条件.
    两个存在式均表示至少有一对\((a,b)\)使\(A(a,b)\)为真，
    同一对个体可以按两种顺序作为见证.
    所以两式成立.
\end{proof}

\begin{remark}
\((\forall x)(\exists y)[A(x,y)]
\neq (\exists y)(\forall x)[A(x,y)]\).
\end{remark}
\begin{example}
    用有限论域说明不同种类的量词一般不能交换.
    \begin{solution}
        取\(D=\{0,1\}\)，
        令\(A(x,y)\)表示\(x=y\).
        每个\(x\)都可取\(y=x\)，
        所以\((\forall x)(\exists y)A(x,y)\)为真.
        但每个固定的\(y\)都有一个不同于它的个体\(x\)，
        所以\((\exists y)(\forall x)A(x,y)\)为假.
        因而这两个公式不逻辑等值.
    \end{solution}
\end{example}

\begin{example}
%@see: 《离散数学》（邓辉文） P130 例4-14
    设\(A\)没有自由的\(y\)，
    \(B\)没有自由的\(x\).
证明：\begin{equation}
	(\forall x)(\forall y)[A(x) \limp B(y)]
	= (\exists x)[A(x)] \limp (\forall y)[B(y)].
\end{equation}
\begin{proof}
显然\begin{align*}
	(\forall x)(\forall y)[A(x) \limp B(y)]
	&= (\forall x)(\forall y)[\neg A(x) \lor B(y)] \\
	&= (\forall x)
		[
			(\forall y)
			[
				\neg A(x) \lor B(y)
			]
		] \\
	&= (\forall x)
		[
			\neg A(x)
			\lor
			(\forall y)[B(y)]
		] \\
	&= (\forall x)[\neg A(x)] \lor (\forall y)[B(y)] \\
	&= \neg(\exists x)[A(x)] \lor (\forall y)[B(y)] \\
	&= (\exists x)[A(x)] \limp (\forall y)[B(y)].
	\qedhere
\end{align*}
\end{proof}
\end{example}

\begin{example}
%@see: 《离散数学》（邓辉文） P131 习题4.4 6.(1)
    设\(A\)没有自由的\(y\)，
    \(B\)没有自由的\(x\).
证明：\begin{equation}
	(\forall x)(\forall y)[A(x) \lor B(y)]
	= (\forall x)[A(x)] \lor (\forall y)[B(y)].
\end{equation}
    \begin{proof}
        这里\(A\)没有自由的\(y\)，
        \(B\)没有自由的\(x\).
        固定解释及其余自由变元赋值.
        若右侧真，
        则\(A\)处处真或\(B\)处处真，
        两种情况都使左侧真.
        若右侧假，
        则存在\(a\)使\(A(a)\)假，
        也存在\(b\)使\(B(b)\)假.
        这对\((a,b)\)使\(A(a)\lor B(b)\)假，
        所以左侧也假.
        故两侧总同真同假.
    \end{proof}
\end{example}

\begin{example}
%@see: 《离散数学》（邓辉文） P131 习题4.4 6.(2)
    设\(A\)没有自由的\(y\)，
    \(B\)没有自由的\(x\).
证明：\begin{equation}
	(\exists x)(\exists y)[A(x) \land B(y)]
	= (\exists x)[A(x)] \land (\exists y)[B(y)].
\end{equation}
    \begin{proof}
        这里\(A\)没有自由的\(y\)，
        \(B\)没有自由的\(x\).
        若左侧真，
        取见证\((a,b)\)，
        它们分别使\(A(a),B(b)\)为真，
        所以右侧真.
        若右侧真，
        分别取两个存在式的见证\(a,b\)，
        就使\(A(a)\land B(b)\)真，
        所以左侧真.
        两个见证不必相同，
        这两个方向对任意解释成立，
        故两式等值.
    \end{proof}
\end{example}

\begin{example}
    %@see: 《离散数学》（邓辉文） P131 习题4.4 6.(3)
    设\(A\)没有自由的\(y\)，
    \(B\)没有自由的\(x\).
    证明下式（全称量词补上指导变元，
    第二个谓词写作\(B(y)\)）：
    \begin{align*}
        (\exists x)(\exists y)[A(x)\limp B(y)] &= (\forall x)A(x)\limp(\exists y)B(y).
    \end{align*}
    \begin{proof}
        利用条件式的析取表示、量词辖域变换和量词转换，
        有
        \begin{align*}
            (\exists x)(\exists y)[A(x)\limp B(y)] &= (\exists x)(\exists y)[\neg A(x)\lor B(y)] \\
            &= (\exists x)[\neg A(x)\lor(\exists y)B(y)] \\
            &= (\exists x)\neg A(x)\lor(\exists y)B(y) \\
            &= \neg(\forall x)A(x)\lor(\exists y)B(y) \\
            &= (\forall x)A(x)\limp(\exists y)B(y).
        \end{align*}
        两次辖域变换均满足自由变元限制，
        并使用了非空论域条件.
    \end{proof}
\end{example}

\section{前束范式}
\subsection{前束范式的概念}
\begin{definition}
    %@see: 《离散数学》（邓辉文） P131 定义4-4
    设\(A\)是谓词公式.
    若\(A=(Q_1x_1)(Q_2x_2)\dotsm(Q_nx_n)[B]\)，
    其中\(n\geq0\)，
    每个\(Q_i\)为\(\forall\)或\(\exists\)，
    \(B\)不含量词，
    则称右侧公式为\(A\)的\DefineConcept{前束范式}（prenex normal form）.
    \(n=0\)时，
    该式就是无量词公式\(B\).
\end{definition}

直观地说，谓词公式的前束范式是将它的所有量词放到最前面去作用整个\(B\).
例如，\(A = (\forall x)[P(x)] \limp Q(x,y)\)不是\(A\)的前束范式.

\subsection{前束范式的计算}
前束范式的计算步骤如下：
\begin{enumerate}
    \item 消去\(\neg,\land,\lor\)以外的联结词.
    \item 使用对合律、德摩根律和量词转换式，
    把否定移到原子公式前面.
    \item 一致地重命名约束变元，
    使不同量词使用不同名字，
    且这些名字不与任何自由变元冲突.
    \item 使用量词辖域的收缩与扩张，
    逐层把量词移到前面，
    每次检查所移量词的变元在另一操作数中不自由出现.
    每个原有量词前缀内部的相对次序必须保留.
\end{enumerate}

\begin{theorem}[前束范式的存在性]
    在非空论域的一阶逻辑中，
    每个公式都有逻辑等值的前束范式，
    而且可保持自由变元不变.
    \begin{proof}
        消去其他联结词后，
        反复使用德摩根律和量词转换式，
        每次把否定推进一层，
        因公式有限，
        最终得到只由文字（原子公式或其否定）、合取、析取和量词组成的公式.
        可用变元名无限而当前公式有限，
        所以总能为每个量词选取互不相同且避开全部自由变元的新名字.
        这种一致改名不改变真值：新、旧量词都逐一遍历同一论域，
        被约束位置取得相同的个体.

        对改名后的公式作结构归纳.
        文字已经是前束范式.
        若子公式已化为前束范式，
        在它前面增加一个量词仍是前束范式.
        若两个子公式分别具有前缀\(Q_1x_1,\dotsc,Q_mx_m\)
        和\(R_1y_1,\dotsc,R_ny_n\)，
        无量词部分分别为\(B,C\)，
        考虑它们的合取或析取.
        全部变元已标准化改名，
        所以每个\(x_i\)都不在另一子公式中自由出现，
        每个\(y_j\)也不在第一子公式中自由出现.
        逐次应用已证的辖域等值式，
        先提出第一前缀再提出第二前缀，
        得到
        \begin{align*}
            (Q_1x_1)\dotsm(Q_mx_m)(R_1y_1)\dotsm(R_ny_n)[B\circ C],
        \end{align*}
        其中\(\circ\)为\(\land\)或\(\lor\).
        这就是前束范式.
        每次变换都不捕获自由变元，
        因而自由变元保持不变.
        归纳完成.
        该构造不许可任意交换\(\forall\)和\(\exists\)，
        只是在各步自由变元条件成立时扩张辖域.
    \end{proof}
\end{theorem}

\begin{example}
    求\((\forall x)P(x)\limp Q(x,y)\)的一个前束范式.
    \begin{solution}
        右侧\(Q(x,y)\)中的\(x\)是自由出现，
        因此不能直接把左侧的\(\forall x\)移到整个公式前面.
        先把左侧约束变元改名为新的\(z\)，
        然后计算
        \begin{align*}
            (\forall x)P(x)\limp Q(x,y) &= (\forall z)P(z)\limp Q(x,y) \\
            &= \neg(\forall z)P(z)\lor Q(x,y) \\
            &= (\exists z)\neg P(z)\lor Q(x,y) \\
            &= (\exists z)[\neg P(z)\lor Q(x,y)].
        \end{align*}
        最后一步合法，
        因为\(z\)在\(Q(x,y)\)中不自由出现.
        所得公式仍以\(x,y\)为自由变元.
    \end{solution}
\end{example}

\section{谓词逻辑中的推理}
\subsection{逻辑蕴含式}
\begin{example}
    %@see: 《离散数学》（邓辉文） P133 例4-19
    设项\(t\)对\(x\)在\(A\)中可自由代入.
    证明\((\forall x)A(x)\implies A(t)\).
    \begin{proof}
        由\cref{example:数理逻辑.全称实例化}，
        \((\forall x)A(x)\limp A(t)\)在所有解释及自由变元赋值下都为真.
        因而任何使前件为真的解释及赋值都使后件为真，
        这正是所需的逻辑蕴含.
    \end{proof}
\end{example}

\begin{example}
%@see: 《离散数学》（邓辉文） P133 例4-20
证明：\((\exists x)(\forall y)[A(x,y)] \implies (\forall y)(\exists x)[A(x,y)]\).
\begin{proof}
任意给定个体域\(D\)上的解释\(I\)，
假定\((\exists x)(\forall y)[A(x,y)]\)在解释\(I\)下取真，
则存在\(d_0 \in D\)，
对于任意\(d \in D\)，
均有\(A(d_0,d)\)取真，
于是\((\forall y)(\exists x)[A(x,y)]\)在解释\(I\)下取真.
\end{proof}
\end{example}

\begin{example}
%@see: 《离散数学》（邓辉文） P133 例4-21
证明：\((\exists y)(\forall x)[A(x,y)]\)不是\((\forall x)(\exists y)[A(x,y)]\)的有效结论.
\begin{proof}
令个体域\(D=\mathbb{R}\)，
\(A(x,y): x>y+3\)，
则\((\forall x)(\exists y)[A(x,y)]\)表示
“对于任意实数\(x\)，均存在实数\(y\)，使得\(x>y+3\)”，
        取\(y=x-4\)即可，
        所以它是真命题；
而\((\exists y)(\forall x)[A(x,y)]\)表示
“存在实数\(y\)，对于任意实数\(x\)，都有\(x>y+3\)”，
        对每个候选\(y\)取\(x=y+3\)即可否定其全称要求，
        所以它是假命题；
所以\((\forall x)(\exists y)[A(x,y)] \limp (\exists y)(\forall x)[A(x,y)]\)不是永真式，
因此\((\exists y)(\forall x)[A(x,y)]\)不是\((\forall x)(\exists y)[A(x,y)]\)的有效结论.
\end{proof}
\end{example}

\subsection{基本推理规则}
命题逻辑中的基本推理规则可以很方便地推广到谓词逻辑.

谓词逻辑有两个非常重要的、与量词有关的逻辑蕴含式.
\begin{theorem}\label{theorem:一阶逻辑.基本推理规则}
%@see: 《离散数学》（邓辉文） P133 定理4-1
设谓词公式\(A,B\)都含自由变元\(x\)，
则\begin{gather}
	(\forall x)[A(x)] \lor (\forall x)[B(x)]
	\implies (\forall x)[A(x) \lor B(x)], \\
	(\exists x)[A(x) \land B(x)]
	\implies (\exists x)[A(x)] \land (\exists x)[B(x)].
\end{gather}
    \begin{proof}
        固定任意解释及其余自由变元赋值.
        对于第一条，
        若前提为真，
        则至少一个全称式为真.
        无论是所有个体满足\(A\)，
        还是所有个体满足\(B\)，
        每个个体均满足\(A\lor B\)，
        所以结论为真.
        对于第二条，
        前提的一个见证同时满足\(A,B\)，
        因而它可以分别作为两个存在式的见证，
        所以结论的合取为真.
        两条规则都把真前提保留为真结论，
        故成立.
    \end{proof}
\end{theorem}

\begin{example}
%@see: 《离散数学》（邓辉文） P134 例4-22(1)
设谓词公式\(A,B\)都含自由变元\(x\)，
举例说明\begin{equation*}
	\neg(\forall x)[A(x)] \implies (\forall x)[\neg A(x)]
\end{equation*}不成立.
\begin{solution}
因为\(
	%\cref{equation:数理逻辑.量词转换1}
	\neg(\forall x)[A(x)]
	= (\exists x)[\neg A(x)]
\)，
而\((\exists x)[\neg A(x)]\)显然不能推出\((\forall x)[\neg A(x)]\).
例如，取个体域\(D = \mathbb{Z}\)，
令\(A(x): \text{$x$是偶数}\)，
则\((\forall x)[A(x)]\)是假命题，
从而\(\neg(\forall x)[A(x)]\)是真命题，
而\((\forall x)[\neg A(x)]\)表示“任意整数都不是偶数”，它是假命题，
因此结论\(\neg(\forall x)[A(x)] \implies (\forall x)[\neg A(x)]\)不成立.
\end{solution}
\end{example}

\begin{example}
%@see: 《离散数学》（邓辉文） P134 例4-22(2)
设谓词公式\(A,B\)都含自由变元\(x\)，
举例说明\begin{equation*}
	(\exists x)[A(x) \limp B(x)] \implies (\exists x)[A(x)] \limp (\exists x)[B(x)]
\end{equation*}不成立.
    \begin{solution}
取个体域\(D = \{1,2\}\)，
解释\(
	\frac{A(1)}1,
	\frac{A(2)}0,
	\frac{B(1)}0,
	\frac{B(2)}0
\).
这时因为\(A(2) \limp B(2)\)是真命题，
于是\((\exists x)[A(x) \limp B(x)]\)是真命题；
而\(A(1)\)为真，
即\((\exists x)[A(x)]\)为真；
但\((\exists x)[B(x)]\)为假；
所以\((\exists x)[A(x)] \limp (\exists x)[B(x)]\)在上述解释下为假，
因此\((\exists x)[A(x) \limp B(x)] \implies (\exists x)[A(x)] \limp (\exists x)[B(x)]\)不成立.
    \end{solution}
\end{example}

\subsection{自然推理系统}
谓词逻辑的自然推理系统是对命题逻辑的自然推理系统的一种推广.
相比于命题逻辑，
谓词逻辑的初始符号，增加了函词、谓词、量词；
新建了谓词公式的形成规则；
依然没有公理；
基本推理规则除了\cref{theorem:一阶逻辑.基本推理规则} 引入的两个逻辑蕴含式以外，
还有以下4个与量词有关的基本推理规则.
记\(\Gamma\)为尚未解除的前提，
\(\Gamma\vdash C\)表示存在由这些前提推出\(C\)的形式推导，
不同于表示语义蕴含的\(\Gamma\models C\).
这里的\(c\)可以是为当前子推导临时引入的新参数，
不是预先指定、具有特殊性质的个体.

\begin{enumerate}
    \item \textbf{全称量词消去（US）}：
    由\((\forall x)A(x)\)可推出\(A(t)\)，
    其中\(t\)对\(x\)在\(A\)中可自由代入.
    \begin{proof}
        全称式为真时，
        \(A\)对每个个体都真，
        特别对项\(t\)表示的个体为真，
        无捕获代入保证得到\(A(t)\).
    \end{proof}

    \item \textbf{全称量词产生（UG）}：
    若\(\Gamma\vdash A(c)\)，
    且\(c\)是没有出现在\(\Gamma\)及原公式\(A\)中的新参数，
    则可得\(\Gamma\vdash(\forall x)A(x)\).
    \begin{proof}
        固定任意使\(\Gamma\)为真的解释及赋值.
        由于\(c\)未出现在这些前提中，
        可将\(c\)解释为论域中的任意个体而不改变前提真值.
        已得子推导保证每种扩充解释都使\(A(c)\)为真，
        所以\(A\)对每个个体都真，
        即全称式为真.
        因而这条规则在子推导可靠时保持可靠性.
    \end{proof}

    \item \textbf{存在量词消去（ES）}：
    已有\(\Gamma\vdash(\exists x)A(x)\)时，
    取未在\(\Gamma,A,C\)中出现的新参数\(c\)，
    在临时假设\(A(c)\)下推出\(C\).
    若\(C\)不含\(c\)，
    则解除该临时假设并得到\(\Gamma\vdash C\).
    \begin{proof}
        固定任意使\(\Gamma\)为真的解释及赋值.
        存在式为真，
        所以有见证\(d\).
        把新参数\(c\)解释为\(d\)，
        不改变\(\Gamma\)的真值，
        且使临时假设\(A(c)\)为真.
        子推导于是使\(C\)为真.
        \(C\)不含\(c\)，
        所以退回原解释后\(C\)仍真，
        故可以解除临时见证假设.
        这不是无条件的逻辑蕴含\((\exists x)A(x)\implies A(c)\).
    \end{proof}

    \item \textbf{存在量词产生（EG）}：
    由\(A(t)\)可推出\((\exists x)A(x)\)，
    其中\(t\)对\(x\)在\(A\)中可自由代入.
    \begin{proof}
        \(A(t)\)为真时，
        项\(t\)表示的个体就是存在式的见证，
        所以结论为真.
    \end{proof}
\end{enumerate}

\begin{example}
    说明为什么不能把UG和ES写成不带条件的普通逻辑蕴含.
    \begin{solution}
        取\(D=\{0,1\}\)，
        令\(A\)仅对\(0\)成立.
        若把已有常量\(c\)解释为\(0\)，
        则\(A(c)\)真而\((\forall x)A(x)\)假，
        所以不能从某个特定个体满足\(A\)就全称推广.
        若把\(c\)解释为\(1\)，
        则\((\exists x)A(x)\)真而\(A(c)\)假，
        所以不能把存在见证任意指定为已有常量.
        新参数条件和临时假设的解除正是为排除这些错误.
    \end{solution}
\end{example}

\begin{example}
%@see: 《离散数学》（邓辉文） P134 例4-23
证明苏格拉底三段论推理的有效性.
\begin{proof}
令\(s: \text{苏格拉底},
P(x): \text{$x$是人},
D(x): \text{$x$是要死的}\)，
则苏格拉底三段论
“所有的人都是要死的，苏格拉底是人，所以，苏格拉底是要死的”
可以转写为\begin{equation*}
	(\forall x)[P(x) \limp D(x)],
	P(s)
	\implies
	D(s).
\end{equation*}
于是
\begin{center}
	\begin{tblr}{cp{4cm}l}
		(1) & \(P(s)\) & P \\
		(2) & \((\forall x)[P(x) \limp D(x)]\) & P \\
		(3) & \(P(s) \limp D(s)\) & US(2) \\
		(4) & \(D(s)\) & T(1)(3)I \\
	\end{tblr}
\end{center}
\end{proof}
\end{example}

\begin{example}
    %@see: 《离散数学》（邓辉文） P134 例4-24
    证明\((\forall x)[F(x)\limp G(x)],(\exists x)F(x)\implies(\exists x)G(x)\).
    \begin{proof}
        取没有出现在前提及目标结论中的新参数\(c\)，
        开启一个以\(F(c)\)为临时假设的子推导.
        对全称前提使用US，
        得\(F(c)\limp G(c)\)；结合\(F(c)\)用假言推理得\(G(c)\)，
        再用EG得\((\exists x)G(x)\).
        这个结论不含\(c\)，
        因而可对前提\((\exists x)F(x)\)使用ES，
        解除临时假设\(F(c)\)，
        得到只依赖两个原前提的结论\((\exists x)G(x)\).
    \end{proof}
\end{example}

\begin{example}
%@see: 《离散数学》（邓辉文） P135 例4-25
证明：\(\neg(\exists x)[F(x) \land H(x)],
(\forall x)[G(x) \limp H(x)]
\implies
(\forall x)[G(x) \limp \neg F(x)]\).
\begin{proof}
用构造法.
        取未在前提及原公式中出现的新参数\(c\)，
        其解释可为论域中的任意个体，
        所以最后一步满足UG的新参数条件.
\begin{center}
	\begin{tblr}{cp{4cm}l}
		(1) & \(\neg(\exists x)[F(x) \land H(x)]\) & P \\
		(2) & \((\forall x)[\neg F(x) \lor \neg H(x)]\) & T(1)E \\
		(3) & \(\neg F(c) \lor \neg H(c)\) & US(2) \\
		(4) & \(H(c) \limp \neg F(c)\) & T(3)E \\
		(5) & \((\forall x)[G(x) \limp H(x)]\) & P \\
		(6) & \(G(c) \limp H(c)\) & US(5) \\
		(7) & \(G(c) \limp \neg F(c)\) & T(4)(6)I \\
		(8) & \((\forall x)[G(x) \limp \neg F(x)]\) & UG(7) \\
	\end{tblr}
\end{center}
\end{proof}
\end{example}
